% macro conflict, first definition kept: \newcommand{\Hom}{\underline{\mathsf{Hom}}}
% macro conflict, first definition kept: \newcommand{\isoto}{\overset{\sim}{\longrightarrow}}
% macro conflict, first definition kept: \newcommand{\Ob}{\mathop{\mathsf{Ob}}}
% macro conflict, first definition kept: \newcommand{\Dia}{\textsf{Dia}}
% macro conflict, first definition kept: \newcommand{\Cat}{\mathcal{C}\mathit{at}}
% macro conflict, first definition kept: \newcommand{\Hot}{\textsf{Hot}}
% macro conflict, first definition kept: \newcommand{\Dia}{\textsf{Dia}}
% macro conflict, first definition kept: \newcommand{\Ob}{\textsf{Ob}}
% macro conflict, first definition kept: \newcommand{\Hom}{\textsf{Hom}}
% macro conflict, first definition kept: \newcommand{\HOT}{\textsf{HOT}}
% macro conflict, first definition kept: \newcommand{\Hot}{\textsf{Hot}}
% macro conflict, first definition kept: \newcommand{\Dia}{\textsf{Dia}}
% macro conflict, first definition kept: \newcommand{\Dia}{\sfup{Dia}}
% macro conflict, first definition kept: \newcommand{\Hot}{\sfup{Hot}}
% macro conflict, first definition kept: \newcommand{\HOT}{\sfup{HOT}}
% macro conflict, first definition kept: \newcommand{\Ob}{\sfup{Ob}}
% macro conflict, first definition kept: \newcommand{\Hom}{\sfup{Hom}}
% macro conflict, first definition kept: \newcommand{\Cat}{$\mathcal{C}$\textit{at}}
% macro conflict, first definition kept: \newcommand{\HOT}{\textsf{HOT}}
% macro conflict, first definition kept: \newcommand{\Hot}{\textsf{Hot}}
% macro conflict, first definition kept: \newcommand{\Dia}{\textsf{Dia}}
\documentclass[11pt,leqno]{article}
\usepackage{fontspec}
\usepackage[english,french]{babel}
\usepackage{amsmath}
\usepackage{amssymb}
\usepackage{amsthm}
\usepackage{mathrsfs}
\usepackage{tikz-cd}
\usepackage{enumitem}
\usepackage[a4paper,left=23mm,right=31mm,top=23mm,bottom=24mm,marginparwidth=22mm,marginparsep=4mm]{geometry}
\usepackage{microtype}
\usepackage{marginnote}
\usepackage{xurl}
\setlength{\emergencystretch}{2em}
\setlength{\parindent}{1.25em}
\setlength{\parskip}{0.15em}
\raggedbottom
\AtBeginDocument{\ifdefined\dcomment\input{DER_I-b_frontmatter.tex}\fi}

\frenchsetup{AutoSpacePunctuation=false}

\newcommand{\printedpage}[1]{\marginnote{\scriptsize\sffamily ed.\ #1}}
\newcommand{\msp}[1]{%
  \begingroup\def\thismsp{#1}\def\blankmsp{page 114 vide}%
  \ifx\thismsp\blankmsp
    \marginnote{\scriptsize\sffamily [114: vide]}[1.3\baselineskip]%
  \else\def\followingmsp{page 115}%
    \ifx\thismsp\followingmsp
      \marginnote{\scriptsize\sffamily [page 115]}[2.6\baselineskip]%
    \else\marginnote{\scriptsize\sffamily [#1]}[1.3\baselineskip]\fi
  \fi\endgroup}
\newcommand{\DD}{\ensuremath{\mathbb{D}}}
\newcommand{\cA}{\ensuremath{\mathcal{A}}}
\newcommand{\Dia}{\ensuremath{\mathsf{Dia}}}
\newcommand{\Ob}{\ensuremath{\mathsf{Ob}}}
\newcommand{\Hom}{\ensuremath{\mathsf{Hom}}}
\newcommand{\hot}{\ensuremath{\mathsf{hot}}}
\newcommand{\Hot}{\ensuremath{\mathsf{Hot}}}
\newcommand{\HOT}{\ensuremath{\mathsf{HOT}}}
\newcommand{\isoto}{\overset{\sim}{\to}}
\newcommand{\isolongto}{\overset{\sim}{\longrightarrow}}
\newcommand{\D}{\ensuremath{\mathbb{D}}}
\newcommand{\Fl}{\ensuremath{\mathsf{Fl}}}
\newcommand{\Cat}{\ensuremath{\mathcal{C}at}}
\newcommand{\bs}{\backslash}
\newcommand{\RR}{\ensuremath{\mathbb{R}}}
\newcommand{\sfup}[1]{\textsf{\upshape #1}}
\begin{document}
% ---- from der1_p036-040.tex
\printedpage{34}
% head: A. Grothendieck : Dérivateurs, Ch. I, Édition des Universités de Montpellier II et Paris VII
% foot: Version 8 juillet 2001
\noindent est à fibres contractiles.

Autre situation standard où le changement de base marche (mais où les foncteurs qui interviennent n'ont aucune raison d'être propres ou lisses).

Soient
\[
f : I \longrightarrow I' \quad , \qquad g : J \longrightarrow J'
\]
dans $\Dia$, d'où le diagramme (comparer page 23)
\[
\begin{tikzcd}
X' = I \times J \arrow[r, "G'"] \arrow[d, "F'"'] & I \times J' = X \arrow[d, "F"] \\
Y' = I' \times J \arrow[r, "G"] & I' \times J' = Y
\end{tikzcd}
\quad .
\]
Ce diagramme est tel que pour tout dérivateur à gauche

\noindent\msp{page 60}

{\raggedright
\noindent on a le changement de base
\[
G^*F_* \simeq F'_*{G'}^* \quad ,
\]
la raison ici étant que les foncteurs induits
\[
X'/y' \longrightarrow X/y
\]
sont des hot-équivalences \underline{\underline{co}}strictes, \textit{i.e.} que les catégories $K$ (page 49) sont contractiles. En effet, si $x_0 = (i_0, j'_0) \in \Ob X$, $y'_0 = (i'_0, j_0) \in \Ob Y'$, $u_0 = (u'_0, v'_0) : F(x_0) \to G(y'_0)$, où $u'_0 : f(i_0) \to i'_0$, $v'_0 : j'_0 \to g(j_0)$ sont donnés, on vérifie que la catégorie $K$ cherchée est isomorphe simplement à $(i_0, u'_0)\backslash(I/i'_0) \times (j_0, v'_0)\backslash(J/j'_0)$ qui est contractile ayant un objet initial \texttt{[plutôt $(i_0, u'_0)\backslash(I/i'_0) \times (j'_0\backslash J)/(j_0, v'_0)$ qui est contractile étant produit d'une catégorie ayant un objet initial et d'une catégorie ayant un objet final]}. Ouf!\par}

Je m'aperçois d'un très gros canular. Je croyais me rappeler que pour $f : X \to Y$ qui est une hot-équivalence stricte, on avait pour tout $F \in \DD(Y)$, $F \isoto f_*f^*F$ (par exemple pour $\DD = \mathsf{Der}(\mathit{Ab})$, ou $\DD = \Hot$), \textit{i.e.} $f^*$ pleinement fidèle, d'où aussi que $\Gamma_{\DD}(Y, F) \to \Gamma_{\DD}(X, f^*F)$ est \texttt{[un]} isomorphisme. Mais ce n'est vrai que si $F$ est un coefficient \textit{constant} sur $Y$, \textit{i.e.} de la forme $q^*(M)$ ($q : Y \to e$, $M \in \DD(e) = \cA_{\DD}$). C'est faux, pour $F$ général, déjà dans le cas de l'immersion $f : \{\varphi\} \to Y$, où $\varphi$ est un objet initial. Alors, $f$ est toujours une hot-équivalence stricte, mais l'homomorphisme $\Gamma(Y, F) \to \Gamma(X, f^*F)$ n'est en général pas un isomorphisme.

\clearpage
\printedpage{35}
% head: A. Grothendieck : Dérivateurs, Ch. I, Édition des Universités de Montpellier II et Paris VII
% foot: Version 8 juillet 2001
\noindent\msp{page 61}

\subsection*{Notions de $\mathbb{D}$-équivalence et de $\mathbb{D}$-contractibilité.}

On va donc reprendre les questions d'acyclicité et de changement de base.

\noindent\textbf{Définition 1.} Soit $\DD$ un prédérivateur, et soit
\[
\begin{tikzcd}
X \arrow[rr, "f"] \arrow[dr, "p"'] & & Y \arrow[dl, "q"] \\
& e &
\end{tikzcd}
\]
dans $\Dia$. On dit que $f$ est une $\DD$-\textit{équivalence} si le foncteur
\[
F \longmapsto f^*F \quad , \qquad \DD(Y) \longrightarrow \DD(X)
\]
induit un foncteur pleinement fidèle sur la sous-catégorie de $\DD(Y)$ formée des $\DD$-coefficients \textit{constants}, \textit{i.e.} si pour tous $M, N \in \Ob \cA = \DD(e)$ (“coefficients absolus”), on a
\[
\Hom(q^*M, q^*N) \isolongto \Hom(p^*M, p^*N) \quad . \tag*{(1)}
\]
Supposons que $p_*$ et $q_*$ existent. On posera, pour un faisceau $F$ sur $X$
\[
H^*_{\DD}(X, F) \overset{\text{déf}}{=} p_*(F) \qquad \text{(et de même } \ H^{\DD}_*(X, F) = p_!(F) \ \text{).}
\]
N.B. Et si $F = p^*M$, $M \in \cA$, on note simplement $H^*_{\DD}(X, M)$, $H^{\DD}_*(X, M)$.

\noindent L'homomorphisme (1) s'interprète par adjonction comme
\[
\begin{tikzcd}[row sep=small]
\Hom(M, q_*q^*N) \arrow[r] \arrow[d, equal] & \Hom(M, p_*p^*N) \arrow[d, equal] \\
\Hom(M, H^*_{\DD}(Y, q^*N)) \arrow[r] & \Hom(M, H^*_{\DD}(X, p^*N))
\end{tikzcd}
\]
déduit de l'homomorphisme
\[
q_*q^*N \longrightarrow p_*p^*N \quad , \qquad \textit{i.e.} \quad H^*_{\DD}(Y, N) \longrightarrow H^*_{\DD}(X, N) \quad . \tag*{(2)}
\]
Donc dire que (1) est un isomorphisme pour tous $M, N$ signifie que (2) est un isomorphisme pour tout $N$. De même, si $p_!$, $q_!$ existent cela équivaut à dire que
\[
p_!p^*N \longrightarrow q_!q^*N \quad , \qquad \textit{i.e.} \quad H^{\DD}_*(X, N) \longrightarrow H^{\DD}_*(Y, N) \tag*{$(2')$}
\]
est un isomorphisme dans $\cA$ pour tout $N \in \cA$.

\clearpage
\printedpage{36}
% head: A. Grothendieck : Dérivateurs, Ch. I, Édition des Universités de Montpellier II et Paris VII
% foot: Version 8 juillet 2001
\noindent\msp{page 62}

\noindent En résumé :

\noindent\textbf{Proposition 1.} \textit{Si $p_*$ et $q_*$ existent, i.e. si $H^*_{\DD}(X, \bullet)$ et $H^*_{\DD}(Y, \bullet)$ existent (il suffit même que ça existe pour des coefficients constants), alors $f$ est une $\DD$-équivalence si et seulement si $f$ induit un isomorphisme pour les $\DD$-cohomologies à coefficients constants de $X$ et $Y$. Si $p_!$ et $q_!$ existent, i.e. si $H^{\DD}_*(X, \bullet)$ et $H^{\DD}_*(Y, \bullet)$ existent, $f$ est une $\DD$-équivalence si et seulement si $f$ induit un isomorphisme pour les $\DD$-homologies à coefficients constants de $X$ et $Y$.}

\noindent\textbf{Remarque 1.} Pour que $f$ soit \texttt{[une]} $\DD$-équivalence, il \textit{suffit} que pour tout $F$ \textit{constant} sur $Y$ le morphisme d'adjonction $F \to f_*f^*F$ soit \texttt{[un]} isomorphisme.

\noindent\textbf{Remarque 2.} Si $f$, $f' : X \to Y$ sont dans une même composante connexe de $\underline{\Hom}(X, Y)$, alors $f$ est $\DD$-équivalence $\Leftrightarrow$ $f'$ $\DD$-équivalence.

\noindent\textbf{Proposition 2.} \textit{Les isomorphismes sont des $\DD$-équivalences. Si $X \overset{f}{\to} Y \overset{g}{\to} Z$, si deux parmis $f$, $g$, $gf$ sont des $\DD$-équivalences, le troisième aussi.} % [sic] parmis -> parmi, spelling

\noindent\textbf{Définition 2.} On dit que $X$ dans $\Dia$ est $\DD$-\textit{asphérique} (ou $\DD$-\textit{contractile}, ou $\DD$-\textit{trivial}) si $X \overset{p}{\to} e$ est une $\DD$-équivalence, \textit{i.e.} si $p^*$ est pleinement fidèle.

Si $p_*$ existe, cela signifie donc que pour $M \in \Ob(\cA)$
\[
M \overset{\mathrm{adj}}{\isolongto} p_*p^*M \quad , \qquad \textit{i.e.} \quad M \isolongto H^*_{\DD}(X, M) \quad ,
\]
et si $p_!$ existe, cela signifie que pour tout $M$
\[
p_!p^*M \overset{\mathrm{adj}}{\isolongto} M \quad , \qquad \textit{i.e.} \quad H^{\DD}_*(X, M) \isolongto M \quad .
\]
\textbf{Corollaire.} \textit{Si $Y$} (resp. $X$) \textit{est contractile, alors $f : X \to Y$ est une $\DD$-équivalence si et seulement si $X$} (resp. $Y$) \textit{est contractile.}

\noindent\textsc{Exemple}. Prenons $\DD = \HOT$. (On peut prendre $\Dia$ = catégorie des petites catégories). Alors, pour $M \in \cA = \Hot$ et $X$ petite catégorie, on a
\[
H^*_{\HOT}(X, M) \simeq \Hom_{\cA}(\hot(X), M) \quad ,
\]
où
\[
\hot : \mathcal{C}\mathit{at} \longrightarrow \cA = \Hot
\]
est le foncteur canonique qui associe à tout $X$ son “type d'homotopie”. On a
\[
H^{\HOT}_*(X, M) = \hot(X) \times M \quad .
\]

\noindent\msp{page 63}

\noindent On en déduit que $f : X \to Y$ est une $\HOT$-équivalence (au sens de la définition 1) si et seulement si c'est une $\hot$-équivalence (ou “équivalence faible”) \textit{i.e.} si et seulement si
\[
\hot(f) : \hot(X) \isolongto \hot(Y)
\]
\printedpage{37}
% head: A. Grothendieck : Dérivateurs, Ch. I, Édition des Universités de Montpellier II et Paris VII
% foot: Version 8 juillet 2001
est un isomorphisme.

\noindent\textbf{Proposition 2.} \textit{Soit $f : X \to Y$. Conditions équivalentes (si $\DD$} \texttt{[est un]} \textit{dérivateur} à gauche). % Doubt resolved: attached PDF leaf 1 (printed p.37) has upright roman à gauche).

a) \textit{$f$ est une $\DD$-équivalence, et le reste après tout changement de base $Y' \to Y$ qui est un isomorphisme local (par exemple pour $Y/F \to Y$, $F \in \Ob \widehat{Y}$).}

b) \textit{Pour tout $y \in \Ob Y$, le foncteur induit $X/y \to Y/y$ est une $\DD$-équivalence.}

c) \textit{Pour tout $y \in \Ob Y$, $X/y$ est $\DD$-contractile.}\footnote[25]{Autres conditions équivalentes :\par
\indent d) \textit{Pour tout faisceau \underline{constant} $M_Y$ sur $Y$, on a $M_Y \isoto f_*f^*M_Y$.}\par
\indent d$'$) \textit{Comme} d) \textit{avec faisceau \underline{localement} constant.}}

\noindent N.B. On dit que $g : Z \to Y$ est un \textit{isomorphisme local} si pour tout $z \in \Ob Z$, posons $y = g(z)$, le morphisme induit sur les catégories localisées en $z$ et $y$
\[
Z/z \isolongto Y/y
\]
est un isomorphisme. Par exemple toute immersion ouverte, tout foncteur induit $Y/y \to Y$ (foncteur de localisation) sont des isomorphismes locaux.

Il en résulte a) $\Rightarrow$ b) trivialement.

b) $\Leftrightarrow$ c). Résulte du corollaire de la page précédente, si l'on sait que $Y/y$ est $\DD$-contractile\footnote[26]{Cf. cor. 1 de \textbf{Der 5} (g)}. Cela résultera de \textbf{Der 5} ci-dessous.

\noindent\msp{page 64}

c) $\Rightarrow$ a) “Passage du local au global”. Comme l'hypothèse c) est stable par changement de base par immersion ouverte, il suffit de prouver que c) implique que $f$ est une $\DD$-équivalence, \textit{i.e.} que pour tout $F = M_Y$ ($M \in \Ob \cA$), coefficient constant sur $Y$, l'homomorphisme d'adjonction
\[
M_Y \longrightarrow f_*f^*M_Y = f_*M_X
\]
est \texttt{[un]} isomorphisme. Il suffit de prouver que c'est un isomorphisme fibre par fibre (\textbf{Der 2}), donc que pour tout $y \in Y$,
\[
M_y \longrightarrow f_*(M_X)_y
\]
est un isomorphisme. Par \textbf{Der 4} (g), on a
\[
f_*(M_X)_y \isolongto \Gamma(X/y, M) \quad ,
\]
et l'hypothèse que $X/y$ est $\DD$-contractile donne ce qu'on veut.

\noindent\textbf{Proposition 2$'$.} \textit{Dualement, si $\DD$ est un dérivateur à gauche} \texttt{[plutôt à droite]}\textit{, conditions équivalentes sur $f : X \to Y$ :}

a) \textit{$f$ $\DD$-équivalence,} \texttt{[et]} \textit{le reste par tout changement de base $g : Y' \to Y$ qui est un “isomorphisme colocal” (i.e. pour tout $y' \in \Ob Y'$, le foncteur induit sur les catégories “colocalisées” $y'\backslash Y' \isolongto y\backslash Y$ (où $y = g(y')$)} \texttt{[est un isomorphisme]}\textit{).}

b) \textit{Pour tout $y \in Y$, le foncteur induit $y\backslash X \to y\backslash Y$ est une $\DD$-équivalence.}

c) \textit{Pour tout $y \in Y$, $y\backslash Y$ est $\DD$-contractile.}

\noindent Pour le prouver, on doit utiliser que les $y\backslash Y$ sont $\DD$-contractiles (cf. cor. 1 de \textbf{Der 5} (d)).

\clearpage
\printedpage{38}
% head: A. Grothendieck : Dérivateurs, Ch. I, Édition des Universités de Montpellier II et Paris VII
% foot: Version 8 juillet 2001
\noindent\msp{page 65}

\noindent\textbf{Définition 3.} Soit $\DD$ un prédérivateur (peut-être même un dérivateur à gauche ou à droite). On dit que $f$ est une $\DD$-équivalence \textit{stricte} si $f$ satisfait la condition a) de \texttt{[la]} proposition 2. On dit que c'est une $\DD$-équivalence \textit{costricte} si $f$ satisfait celle de \texttt{[la]} proposition $2'$; on dit aussi parfois que $f$ est $\DD$-cofinal. (La condition duale, \textit{i.e.} $\DD$-équivalence stricte, s'appelerait co-$\DD$-cofinal!)\footnote[27]{Meilleure terminologie : $f$ est $\DD$-\textit{asphérique} ($\DD$-équivalence stricte), $\DD$-\textit{coasphérique} ($\DD$-équivalence co-stricte).} % [sic] s'appelerait -> s'appellerait, spelling

\noindent\textbf{Proposition 3.} \textit{Soit $X \overset{f}{\to} Y \overset{g}{\to} Z$. Si $f, g$ sont des $\DD$-équivalences strictes, alors $gf$ aussi. Si $gf$ et $f$ sont des $\DD$-équivalences strictes, alors $g$ aussi. Si $\DD$} \texttt{[est un]} \textit{dérivateur à gauche ou à droite, toute équivalence} \texttt{[de catégories]} \textit{est une $\DD$-équivalence stricte.} (Car elle a un adjoint à droite.)

On utilise le fait que la notion d'isomorphisme local est stable par changement de base.

\noindent\msp{page 66}

\noindent\textbf{Proposition 4.} \textit{Soient $\DD$ un dérivateur à droite, $f : X \to Y$ un foncteur dans} $\Dia$\textit{, tel que pour $y \in Y$, $y\backslash X$ soit $\DD$-contractile (donc $f$ une $\DD$-équivalence costricte, sous réserve que $\DD$ soit un dérivateur à droite \ldots\ ). Alors pour tout $F$ dans $\DD(Y)$ (pas nécessairement à coefficients constant) on a} % [sic] constant -> constants, agreement
\[
H^*_{\DD}(Y, F) \isolongto H^*_{\DD}(X, f^*(F))
\]
\textit{est un isomorphisme. Réciproquement, si la condition est vraie pour tout $F$, alors les $y\backslash X$ sont $\DD$-contractiles.}

\noindent\textsc{Démonstration}. Il revient au même de dire que pour tout $M \in \Ob \cA = \DD(e)$, on a
\[
\Hom(M, H^*_{\DD}(Y, F)) \overset{\mathrm{iso}}{\longrightarrow} \Hom(M, H^*_{\DD}(X, f^*(F))) \ .
\]
Mais par adjonction, cette flèche s'interprète comme
\[
\Hom(M_Y, F) \longrightarrow \Hom(M_X, f^*F) \tag*{$(*)$}
\]
D'autre part, si $f_!$ existe, on a
\[
\Hom(M_X, f^*F) = \Hom(f_!(M_X), F) \ ,
\]
et l'homomorphisme $(*)$ est déduit de l'homomorphisme d'adjonction
\[
f_!(M_X) = f_!(f^*(M_Y)) \longrightarrow M_Y \ , \tag*{$(**)$}
\]
en prenant les $\Hom$ dans $F$. Donc la condition équivaut à dire que cet homomorphisme est \texttt{[un]} isomorphisme, quel que soit le coefficient constant $M_Y$ sur $Y$. Pour voir si c'est un isomorphisme, on regarde fibre par fibre (\textbf{Der 2}), or les fibres s'interprètent comme $H^{\DD}_*(y\backslash X, M_X)$ (par \textbf{Der 4} (d)) et comme $M$.

% ---- from der1_p041-045.tex
\printedpage{39}
% head: A. Grothendieck : Dérivateurs, Ch. I, Édition des Universités de Montpellier II et Paris VII
% foot: Version 8 juillet 2001
\msp{page 67}

\noindent Donc la flèche envisagée $(**)$ s'interprète comme
\[
H_*^{\D}(y\backslash X, M_X) \longrightarrow M
\]
et c'est l'homomorphisme d'adjonction $p_!(p^*(M)) \to M$.
\[
\begin{tikzcd}
X \arrow[rr, "f"] \arrow[dr, "p"{swap}] & & Y \arrow[dl, "q"] \\
& e &
\end{tikzcd}
\]

\noindent Donc c'est un isomorphisme pour tout $M$ si et seulement si $y\backslash X$ est $\D$-contractile.

\noindent On va énoncer ce qu'on a prouvé.

\noindent\textbf{Théorème 1.} \textit{Soit}
\[
\begin{tikzcd}
X \arrow[rr, "f"] \arrow[dr, "p"{swap}] & & Y \arrow[dl, "q"] \\
& e &
\end{tikzcd}
\]
\textit{dans} $\Dia$. \textit{On suppose que} $\D$ \textit{est un prédérivateur satisfaisant} \textbf{Der 2}, \textit{que} $p_*, q_*$ \textit{et} $f_!$ \textit{existent, que} $f_!$ \textit{satisfait à} \textbf{Der 4} (d) \textit{pour le calcul des fibres. Conditions équivalentes :}

a) $f$ \textit{est une} $\D$\textit{-équivalence costricte, i.e.} (prop. 2) \textit{pour tout} $y \in \Ob\, Y$, $y\backslash X$ \textit{est} $\D$\textit{-contractile ;}

b) \textit{pour tout} $F \in \Ob\, \D(Y)$, \textit{l'homomorphisme suivant est} \texttt{[un]} \textit{isomorphisme}
\[
H^*_{\D}(Y, F) \isoto H^*_{\D}(X, f^*(F)) \ ;
\]

c) \textit{pour tout} $M \in \Ob(\mathcal{A} = \D(e))$, \textit{l'homomorphisme d'adjonction est} \texttt{[un]} \textit{isomorphisme}
\[
f_!(M_X) = f_!(f^*(M_Y)) \isoto M_Y \ .
\]

Je vais renuméroter les axiomes sur le prédérivateurs. % [sic] le prédérivateurs -> les prédérivateurs, article singulier devant un pluriel

\noindent\textbf{Der 1}, \textbf{Der 2} comme avant (p. 45).

\noindent\textbf{Der 3}\footnote[28]{Axiome des catégories d'indices asphériques, \textit{cf.} variantes \textbf{Der 3} a) b) c) page 97, et surtout pages 111-112.} : Soit $X$ dans $\Dia$, $p : X \to e$. Supposons $X$ contractile.\footnote[29]{N.B. Il suffit de supposer que $X$ a un objet final. Mais ce n'est pas un axiome autodual.} Alors
\[
p^* : M \longmapsto M_X \quad , \qquad p^* : \mathcal{A} = \D(e) \longrightarrow \D(X)
\]
est pleinement fidèle.

\printedpage{40}
% head: A. Grothendieck : Dérivateurs, Ch. I, Édition des Universités de Montpellier II et Paris VII
% foot: Version 8 juillet 2001
\msp{page 68}

\noindent\textbf{Der 4} (g) : existence des $f_*$.

\noindent\textbf{Der 4} (d) : existence des $f_!$.

\noindent\textbf{Der 5} (g) Calcul des fibres de $f_*(F)$.

\noindent\textbf{Der 5} (d) Calcul des fibres de $f_!(F)$.

\noindent L'ancien axiome \textbf{Der 5} (g) résulte de \textbf{Der 2}, \textbf{Der 3}, \textbf{Der 4} (g,d), \textbf{Der 5} (d), grâce au théorème précédent.

Un aspect intéressant du théorème précédent, si on l'envisage comme énonçant l'équivalence de a) et b), c'est que cet énoncé ne fait pas appel à des notions homologiques, mais que sa démonstration y fait appel, sous une forme relativement forte (calcul des fibres de $f_!(F)$). Pour les théorèmes de changement de base pour la $\D$-cohomologie relative, \textit{i.e.} pour les $f_*^{\D}$, on a besoin aussi constamment du calcul des fibres des $f_*(F)$. Donc il semble qu'on ne peut se passer d'aucun des axiomes précédents, y compris \textbf{Der 4} (d) et \textbf{Der 5} (d), même si on ne veut s'intéresser qu'au formalisme des $f^*, f_*$ (aspect cohomologique). Par la suite, on va donc supposer ces axiomes satisfaits.

\noindent\textbf{Théorème 2.} \textit{Soit} $f : X \to Y$ \textit{dans} $\Dia$. \textit{Conditions équivalentes :}

a) \textit{Pour tout} $y_0 \in \Ob\, Y$, \textit{le foncteur canonique}
\[
X_{y_0} \overset{\varphi}{\longrightarrow} X/y_0 \quad (\simeq X \times_Y Y/y_0) \ ,
\]
\textit{où} $X_{y_0}$ \textit{désigne la fibre en} $y_0$, \textit{est une} $\D$\textit{-équivalence co-stricte.}

b) \textit{Pour tout} $x_0 \in \Ob\, X$, \textit{le foncteur induit sur les catégories colocalisées en} $x_0$ \textit{et en} $y_0 = f(x_0)$
\[
x_0\backslash X \longrightarrow y_0\backslash Y
\]
\textit{est à fibres} $\D$\textit{-contractiles.}

\msp{page 69}

c) \textit{La formation des} $f_*F$ \textit{(}$F$ \textit{dans} $\D(X)$\textit{) commute à tout changement de base} $Y' \to Y$, \textit{et ça reste vrai après tout changement de base -- i.e. si on a} $Y'' \to Y' \to Y$, \textit{d'où}
\[
\begin{tikzcd}
X'' \arrow[r] \arrow[d, "f''"{swap}] & X' \arrow[r] \arrow[d, "f'"] & X \arrow[d, "f"] \\
Y'' \arrow[r, "p'"{swap}] & Y' \arrow[r, "p"{swap}] & Y
\end{tikzcd}
\]
\textit{le “théorème de changement de base” est vrai pour} $f'_*$ \textit{et} $p'^{*}$.

c$'$) \textit{La formation des} $f_*(F)$ \textit{(}$F \in \Ob\, \D(X)$\textit{) commute au passage aux fibres et ceci reste vrai après tout changement de base} $p : Y' \to Y$.

c$''$) \textit{Comme} c$'$) \textit{avec} $Y' \to Y$ \textit{un morphisme de localisation} $Y/y \to Y$.

d) \textit{Si} $Y'' \overset{p}{\to} Y' \to Y$, \textit{avec} $p$ \textit{coasphérique, alors} $q : X'' = X \times_Y Y'' \to X' = X \times_Y Y'$ \textit{est coasphérique.}

\printedpage{41}
% head: A. Grothendieck : Dérivateurs, Ch. I, Édition des Universités de Montpellier II et Paris VII
% foot: Version 8 juillet 2001

d$'$) \textit{Comme} d) \textit{avec} $Y'' = e$ (donc $Y'$ est catégorie avec objet final $e$).

d$''$) \textit{Comme} d) \textit{avec} $Y'' = \{1\}$, $Y' = \boldsymbol{\Delta}_1$.

.

\noindent\textsc{Démonstration}. a) $\Leftrightarrow$ b) cf. proposition page 53. Montrons que ces conditions impliquent c). Comme elles sont stables par changement de base (\textit{cf.} démonstration immédiate page 55) donc par \textbf{Der 2} on est ramené à prouver que $H^*_{\D}(X/y, F) \overset{\sim}{\to} H^*_{\D}(Y_y, F)$ $(F \in \Ob\, \D(Y))$, ce qui résulte du théorème 1.\footnote[30]{Donc on utilise \textbf{Der 5} (d).} (Cf. démonstration du théorème page 56, qui marche avec des hypothèses plus générales.)

On a c) $\Leftrightarrow$ c$'$) $\Rightarrow$ c$''$) (l'équivalence c) $\Leftrightarrow$ c$'$) par \textbf{Der 2}).

Prouvons c$''$) $\Rightarrow$ a).

\msp{page 70}

\noindent Considérons le diagramme à carrés cartésiens
\[
\begin{tikzcd}
X_{y_0} \arrow[r, "q"] \arrow[d, "p"{swap}] & X/y_0 \arrow[r] \arrow[d, "f_0"] & X \arrow[d, "f"] \\
\{y_0\} \arrow[r] & Y/y_0 \arrow[r] & Y
\end{tikzcd}
\qquad\qquad .
\]

\noindent L'hypothèse c$''$) prouve que pour tout $F$ dans $\D(X/y_0)$, on a $H^*_{\D}(X/y_0, F) \overset{\sim}{\to} H^*_{\D}(X_{y_0}, F)$. Par le théorème 1, ceci implique que $f$ est un $\D$-équivalence co-stricte. % [sic] un D-équivalence -> une D-équivalence, accord du genre

\noindent\textbf{Corollaire 1.} Énoncé dual du théorème, pour les foncteurs $X_{y_0} \to y_0\backslash X$, $X/x_0 \to Y/y_0$, et le changement de base des $j_!(F)$ \ldots

\noindent\textbf{Définition 4.} Sous les conditions équivalentes du théorème 1, on dit que $f$ est $\D$-\textit{propre.} La condition duale s'appelle $f$ $\D$-\textit{lisse.}

\noindent\textbf{Corollaire 2.} \textit{Soit} $f : X \to Y$ \textit{dans} $\Dia$. \textit{Si} $f$ \textit{est} $\D$\textit{-propre ou} $\D$\textit{-lisse, et à fibres} $\D$\textit{-contractiles, alors} $f$ \textit{est une} $\D$\textit{-équivalence stricte et costricte. Mieux,} $f$ \textit{est une} $\D$\textit{-équivalence} universelle \textit{i.e. est une} $\D$\textit{-équivalence et le reste après tout changement de base.}\footnote[31]{Se généralise : \textbf{Corollaire 2$'$.} \textit{Soit}
\[
\begin{tikzcd}[ampersand replacement=\&]
X \arrow[rr, "f"] \arrow[dr] \& \& Y \arrow[dl] \\
\& S \&
\end{tikzcd}
\]
\textit{avec} $X, Y$, $\D$\textit{-propres sur} $S$. \textit{Si} $f$ \textit{induit une} $\D$\textit{-équivalence sur les fibres c'est une} $\D$\textit{-équivalence. Itou dans le cas lisse.}\par
C'est la proposition 6 ci-dessous.}

\printedpage{42}
% head: A. Grothendieck : Dérivateurs, Ch. I, Édition des Universités de Montpellier II et Paris VII
% foot: Version 8 juillet 2001
\msp{page 71}

\noindent Il suffit de prouver que $f$ est une $\D$-équivalence. Supposons par exemple $f$ $\D$-propre. Il suffit de prouver que si $Y$ est $\D$-contractile, alors $X$ l'est (quitte à l'appliquer aux $Y/y_0$ après changement de base, pour montrer que les $X/y_0$ sont contractiles). Pour ceci, il suffit de prouver que $f$ est une $\D$-équivalence, et on sait que ceci est le cas si pour tout $F$ \textit{constant} sur $Y$, l'homomorphisme $F \overset{\text{adj}}{\longrightarrow} f_*f^*F$ est \texttt{[un]} isomorphisme. En fait, c'est vrai pour \textit{tout} $F$.

\noindent\textbf{Corollaire 3.} \textit{Si} $f : X \to Y$ \textit{est} $\D$\textit{-propre ou} $\D$\textit{-lisse, à fibres} $\D$\textit{-contractiles, alors} $f^* : \D(Y) \to \D(X)$ \textit{est pleinement fidèle, i.e. les homomorphismes d'adjonction}
\[
F \longrightarrow f_*f^*F \quad , \qquad f_!f^*F \longrightarrow F
\]
\textit{dans} $\D(Y)$ \textit{sont des isomorphismes.}

\noindent Il suffit de le voir pour $f$ propre, et cela résulte du calcul des fibres de $f_*f^*F$ \ldots

\noindent\textbf{Proposition 5.} \textit{La notion de foncteur} $\D$\textit{-propre} (resp. $\D$\textit{-lisse}) \textit{est stable par changement de base,} (trivial), \textit{et par composition} (utiliser \texttt{[le]} corollaire 2). \textit{Toute immersion fermée est} $\D$\textit{-propre, toute immersion ouverte} $\D$\textit{-lisse. Plus généralement, tout isomorphisme local est} $\D$\textit{-lisse, tout co-isomorphisme local est} $\D$\textit{-propre.}

\noindent\textbf{Proposition 6.} \textit{Soient} $X, Y$ \textit{sur} $S$, $f : X \to Y$ \textit{un} $S$\textit{-morphisme. Supposons} $X, Y$ \textit{propres sur} $S$, \textit{ou} $X$ \textit{et} $Y$ \textit{lisses sur} $S$, \textit{et que pour tout} $s \in \Ob\, S$, $f_s : X_s \to Y_s$ \textit{soit une} $\D$\textit{-équivalence, alors} $f$ \textit{aussi.}\footnote[32]{\textsc{Démonstration}. On prouvera même que sur la sous-catégorie des coefficients sur $Y$ qui proviennent de $\D(Y)$ \texttt{[plutôt de }$\D(S)$\texttt{]}, $f^* : \D(Y) \to \D(X)$ induit un foncteur pleinement fidèle. Si $f$ \texttt{[est]} propre, je l'interprète comme $q_*q^*F \overset{\sim}{\to} p_*p^*F$ pour tout $F$, et cela résulte du calcul de $p_*$, $q_*$ par passage aux fibres (propre), et de l'hypothèse sur les $f_s$. OK.}

\msp{page 72}

\noindent\textbf{Théorème 3.} \textit{Soit}
\[
\begin{tikzcd}
X \arrow[d, "f"{swap}] & X' \arrow[l, "q"{swap}] \arrow[d, "f'"] \\
Y & Y' \arrow[l, "p"{swap}]
\end{tikzcd}
\]
\textit{cartésien dans} $\Dia$. \textit{Alors le morphisme de changement de base}
\[
p^*f_* \longrightarrow f'_*q^*
\]
\textit{est un isomorphisme si} $f$ \textit{est} $\D$\textit{-propre, ou} $p$ $\D$\textit{-lisse.}

\noindent Le cas $\D$-propre déjà prouvé. Le cas $\D$-lisse : la démonstration p. 56 est OK.

\printedpage{43}
% head: A. Grothendieck : Dérivateurs, Ch. I, Édition des Universités de Montpellier II et Paris VII
% foot: Version 8 juillet 2001

\subsection*{Comparaison de $X$ et $X^{\circ}$.}

Soit $X$ \texttt{[une]} catégorie, on définit une catégorie $\Phi(X)$ et des foncteurs
\[
\begin{tikzcd}[column sep=small]
& & \Phi(X) \arrow[ddl, "p_X"{swap, name=P}] \arrow[ddr, "q_X"{name=Q}] & & \\
|[alias=L]| \text{\footnotesize\begin{tabular}{c}$\Cat$-cofibrant\\(donc propre)\end{tabular}} & & & & |[alias=R]| \text{\footnotesize\begin{tabular}{c}$\Cat$-fibrant\\(donc lisse)\\ \texttt{[plutôt }$\Cat$\texttt{-cofibrant}\\ \texttt{(donc propre)]}\end{tabular}} \\
& X & & X^{\circ} &
\arrow[from=P, to=L]
\arrow[from=Q, to=R]
\end{tikzcd}
\]
qui seront des $\hot$-équivalences, et des $\D$-équivalences si $\D$ est un dérivateur dont le domaine contient $X$ et $\Phi(X)$. On pose
\[
\Ob\, \Phi(X) = \Fl(X) \quad ,
\]
donc mêmes objets que $\mathcal{F}\ell(X) = \underline{\Hom}(\boldsymbol{\Delta}_1, X)$, mais les morphismes sont différents :

\msp{page 73}
\[
\begin{tikzcd}
Y_1 \arrow[r, "u"] & Y_0 \arrow[d, "v_0"] \\
Y'_1 \arrow[u, "v_1"] \arrow[r, "u'"{swap}] & Y'_0
\end{tikzcd}
\]
un morphisme de $u$ dans $u'$ est un couple $v_0 : Y_0 \to Y'_0$, $v_1 : Y'_1 \to Y_1$ tel que $v_0uv_1 = u'$. Ainsi source et but $s(u)$ et $b(u)$ sont l'une contravariante, l'autre covariant en $u$. On prend pour $q$ le foncteur source, pour $p$ le foncteur but. % [sic] l'autre covariant -> l'autre covariante, accord avec « l'une contravariante »

Si on applique ceci à $X^{\circ}$, on trouve
\[
\begin{tikzcd}
& \Phi(X^{\circ}) \arrow[dl, "p_{X^{\circ}}"{swap}] \arrow[dr, "q_{X^{\circ}}"] & \\
X^{\circ} & & (X^{\circ})^{\circ} = X
\end{tikzcd}
\]
On a $\Ob\, \Phi(X^{\circ}) = \Fl(X^{\circ}) \simeq \Fl(X)$
\[
\begin{tikzcd}
Y_1 \arrow[r, "u"] & Y_0 \arrow[d, "v_1"] & \text{dans}\ \ X \\
Y'_1 \arrow[u, "v_0"] \arrow[r, "u'"{swap}] & Y'_0 & \text{dans}\ \ X
\end{tikzcd}
\]
et mêmes homomorphismes au sens de $\Phi(X)$ et au sens de $\Phi(X^{\circ})$. Donc
\[
\left\{
\begin{aligned}
&\Phi(X^{\circ}) = \Phi(X) \\
&p_{X^{\circ}} = q_X \\
&q_{X^{\circ}} = p_X \quad ,
\end{aligned}
\right.
\]

% ---- from der1_p046-050.tex
\printedpage{44}
% head: A. Grothendieck : Dérivateurs, Ch. I, Édition des Universités de Montpellier II et Paris VII
le passage de $X$ à $X^\circ$ revient à échanger $p_X$ et $q_X$. La fibre de $\Phi(X)$ en $x \in \Ob X$ est $\simeq (X/x)^\circ$, celle en $y^\circ \in \Ob X^\circ$ est $\simeq y\bs X$
\[
\left\{
\begin{aligned}
&p_X^{-1}(\{x\}) \simeq (X/x)^\circ\\
&q_X^{-1}(\{y^\circ\}) \simeq y\bs X
\end{aligned}
\right.
\]

{\raggedright
N.B. : $p_X$ est $W$-lisse \texttt{[plutôt $W$-propre]} à fibres $W$-asphériques, donc dans $W^u$ \texttt{[$W^u =$ ensemble des morphismes qui sont dans $W$ et le restent après tout changement de base]}, $q_X$ est $W$-propre à fibres $W$-asphériques, donc dans $W^u$.\par}

\msp{page 74}

On vérifie
\[
\left\{
\begin{aligned}
&\Phi(X)/x \simeq \Phi(X/x) &\qquad& x \in \Ob X\\
&\Phi(X)/x^\circ \simeq \Phi(x\bs X) &\qquad& x^\circ \in \Ob X^\circ
\end{aligned}
\right.
\]

Je dis que $p_X$, $q_X$ sont des $\D$-\textit{équivalences strictes} et \textit{costrictes.} Par les formules précédentes, cela signifie que les $\Phi(X/x)$, $\Phi(x\bs X)$ sont $\D$-contractiles. Donc que si $Z$ a un objet final, $\Phi(Z)$ est $\D$-contractile. Mais considérons
\[
\begin{tikzcd}
\Phi(Z) \arrow[dr, "q_Z"] & \\
& Z^\circ
\end{tikzcd}
\quad ;
\]
il suffit de prouver que c'est une $\D$-équivalence stricte (puisque $Z^\circ$, ayant un objet initial, est $\D$-contractil), donc que les $z_0\bs\Phi(Z)$ sont $\D$-contractiles. On a par les formules précédentes % [sic] contractil -> contractile, adjectif ailleurs écrit « contractile »
\[
z_0\bs\Phi(Z) \simeq \Phi(z_0\bs Z) = \Phi(T) \quad ,
\]
où cette fois $T$ a un objet initial $\emptyset$ et un objet final $e$. Considérons alors l'objet
\[
\emptyset \longrightarrow e
\]
de $\Phi(T)$, c'est un objet final, car pour tout $x \to y$ dans $\Phi(T)$, il y a une unique flèche
\[
\begin{tikzcd}
x \arrow[r] & y \arrow[d] \\
\emptyset \arrow[u] \arrow[r] & e
\end{tikzcd}
\]
% foot: Version 8 juillet 2001

\printedpage{45}
% head: A. Grothendieck : Dérivateurs, Ch. I, Édition des Universités de Montpellier II et Paris VII
\msp{page 75}

Pour résumer

\textbf{Proposition 1 :} \textit{$p_X$ et $q_X$ sont des} \textsf{hot}\textit{-équivalences strictes,} i.e. \textsf{Hot}\textit{-asphériques (donc aussi des $\D$-équivalences strictes si $\Phi(X)$ est dans le domaine de $\D$) à fibres les $(X/x)^\circ$ et $y\bs X$ respectivement (donc à fibres contractiles).}

\textbf{Corollaire 1.} \textit{Pour que $X$ soit $\D$-contractile, il faut et il suffit que $X^\circ$ le soit.}

\hspace*{\parindent}Plus généralement :

\textbf{Corollaire 2.} \textit{Soit $f : X \to Y$. Pour que $f$ soit une $\D$-équivalence, il faut et il suffit que $f^\circ$ le soit.}

\hspace*{\parindent}On utilisera que le diagramme
\[
\begin{tikzcd}
& \Phi(X) \arrow[dl, "p_X"'] \arrow[dr, "q_X"] & \\
X & & X^\circ
\end{tikzcd}
\]
est fonctoriel en $X$. On a donc
\[
\begin{tikzcd}
& \Phi(X) \arrow[rr, "\Phi(f)"] \arrow[dl, "p_X"'] & & \Phi(Y) \arrow[dl, "p_Y"] \arrow[dd, "q_Y"] \\
X \arrow[rr, "f" near end] \arrow[from=1-2, to=3-2, "q_X"' near end, crossing over] & & Y & \\
& X^\circ \arrow[rr, "f^\circ"'] & & Y^\circ
\end{tikzcd}
\]
les deux carrés commutatifs. Comme $p_X$ et $p_Y$ sont des $\D$-équivalences, $f$ est une $\D$-équivalence si et seulement si $\Phi(f)$ l'est, et de même pour $f^\circ$. Ces raisonnements supposent que si $X$ est dans le domaine de $\D$, $\Phi(X)$ aussi. C'est le cas si \Dia{} est formé des catégories ordonnées finies.
% foot: Version 8 juillet 2001

\printedpage{46}
% head: A. Grothendieck : Dérivateurs, Ch. I, Édition des Universités de Montpellier II et Paris VII
\msp{page 76}

\subsection*{$\D$-fibrations}

\textbf{Proposition 1.} \textit{Soit $f : X \to Y$ dans} \Dia\textit{, conditions équivalentes :}

\hspace*{\parindent}a) \textit{Pour tout $y \in Y$, l'inclusion $X_y \to X/y$ est une $\D$-équivalence, et pour toute flèche $y_1 \to y_0$ de $Y$, $X/y_1 \to X/y_0$ est une $\D$-équivalence.}

\hspace*{\parindent}b) \textit{Pour tout coefficient} constant \textit{$F$ sur $Y$, $f_*f^*F$ est localement constant,}\footnote[33]{Par là on entend que pour tout $u : y \to y'$ dans $Y$, $G_{y'} \to G_y$ est \texttt{[un]} isomorphisme. Sans axiomes supplémentaires sur $\D$, il n'est pas clair que cela implique que les ``restrictions'' de $G$ aux $Y/y$ ou $y\bs Y$, soient isomorphes à des coefficients constants. Voir cependant l'axiome \textbf{Der 3} a$'$), page 97, qui assure que c'est pareil.} \textit{et ``se calcule fibre par fibre''} i.e. \textit{l'homomorphisme canonique}
\[
f_*f^*(F)_y \overset{\sim}{\longrightarrow} H^*_{\D}(X_y, F|X_y)
\]
\texttt{[est un isomorphisme]}.

Ces conditions sont équivalentes, et impliquent les conditions de (b) pour tout $F$ \textit{localement} constant sur $Y$.

La première condition dans (a) signifie que pour tout $y \in Y$, et tout $M \in \Ob \mathcal{A}$,
\[
H^*_{\D}(X/y, M) \longrightarrow H^*_{\D}(X_y, M)
\]
est un isomorphisme. Or cette flèche s'interprète comme
\[
f_*(M_X)_y \longrightarrow H^*_{\D}(X_y, M)
\]
est un isomorphisme, \textit{i.e.} que $f_*(M_X)$ se calcule fibre par fibre. Comme $M_X \simeq f^*(M_Y)$, c'est la deuxième condition (b). La deuxième condition dans (a), pour $y_1 \to y_0$ signifie que pour $M$ dans $\mathcal{A}$,
\[
H^*_{\D}(X/y_0, M) \longrightarrow H^*_{\D}(X/y_1, M)
\]
est \texttt{[un]} isomorphisme, or cette flèche s'interprète comme
\[
f_*(M_X)_{y_0} \longrightarrow f_*(M_X)_{y_1} \quad ,
\]
que ce soit toujours un isomorphisme, c'est la première condition (b). Que (b) soit valable aussi pour $F$ localement constant sur $Y$ est immédiat, car c'est une condition de nature locale par rapport à $Y$.

\msp{page 77}

La proposition duale fait appel aux foncteurs
\[
y_0\bs X \longrightarrow y_1\bs X \qquad \text{et} \qquad X_y \longrightarrow y\bs X \quad ,
\]
et aux faisceaux $f_!f^*M$, pour $M$ constant, ou localement constant sur $Y$.
% foot: Version 8 juillet 2001

\printedpage{47}
% head: A. Grothendieck : Dérivateurs, Ch. I, Édition des Universités de Montpellier II et Paris VII
\textbf{Définition 1.} Si les conditions équivalentes de la proposition 1 sont satisfaites, on dit que $f$ est une $\D$-\textit{fibration cohomologique} (ou $\D$-\textit{fibration gauche}). Si les conditions duales sont satisfaites, on dit que c'est une $\D$-\textit{fibration homologique} (ou $\D$-\textit{fibration droite}).

N.B. Les $\D$-fibrations cohomologiques sont stables par tout changement de base qui est un isomorphisme local. Les $\D$-fibrations homologiques par tout changement de base qui est un co-isomorphisme local. Mais ce ne sont pas des notions stables par \textit{tout} changement de base\footnote[34]{? donner exemple!}. Mais la notion de $\D$-fibration cohomologique est stable par changement de base $\D$-lisse, et celle de $\D$-fibration homologique par changement de base $\D$-propre.

\hspace*{\parindent}Je vais essayer de construire un exemple où le changement de base canule, en posant $Y = \boldsymbol{\Delta}_1$. Se donner $X$ sur $\boldsymbol{\Delta}_1$ revient à se donner une sous-catégorie ouverte $X_0$, et son complémentaire fermé $X_1$
\[
\begin{tikzcd}
& X & \\
X_1 \arrow[ur, hook] \arrow[d, dash, dashed] & & X_0 \arrow[ul, hook'] \arrow[d, dash, dashed] \\
1 & & 0 \arrow[ll]
\end{tikzcd}
\quad .
\]
Dire que $X \to Y = \boldsymbol{\Delta}_1$ est une $\D$-fibration cohomologique signifie simplement
\msp{page 78}
que les restrictions
\[
X_0 \overset{j}{\longrightarrow} X \quad , \qquad X_1 \overset{i}{\longrightarrow} X
\]
sont des $\D$-équivalences. Prenons un changement de base $Y' \to Y = \boldsymbol{\Delta}_1$, \textit{i.e.} un $Y'$ avec $Y'_0, Y'_1$
\[
\begin{tikzcd}
& Y' & \\
Y'_1 \arrow[ur, hook] \arrow[d, dash, dashed] & & Y'_0 \arrow[ul, hook'] \arrow[d, dash, dashed] \\
1 & & 0 \arrow[ll]
\end{tikzcd}
\quad .
\]
On voit que la condition $Y'$ lisse sur $\boldsymbol{\Delta}_1$ signifie que pour tout $y'_1 \in \Ob Y'_1$, la catégorie $Y'_0/y'_1$ des objets $y_0$ de $Y'_0$ au dessus de $y'_1$ (\textit{i.e.} munis de $y_0 \to y'_1$) est contractile. On voit que si $Y'_1 \neq \emptyset$ alors $Y'_0 \neq \emptyset$ -- mais un exemple avec $Y'_0 = \emptyset$, où $Y'_1 = \emptyset$, n'est pas intéressant
% foot: Version 8 juillet 2001
\printedpage{48}
% head: A. Grothendieck : Dérivateurs, Ch. I, Édition des Universités de Montpellier II et Paris VII
(car une projection $I \times J \to I$ est toujours une $\D$-fibration cohomologique). Prenons par exemple pour $Y'$
\[
\begin{tikzcd}
& a \arrow[dl] \\
b & \\
& a' \arrow[ul] \\
\underbrace{\qquad}_{\displaystyle Y'_1} & \underbrace{\qquad}_{\displaystyle Y'_0}
\end{tikzcd}
\quad ,
\]
somme amalgamée de deux copies de $Y = \boldsymbol{\Delta}_1$, donc $X' = X \times_Y Y'$ peut se voir comme somme amalgamée $X \amalg_{X_1} X$. Mais ça semble \textit{pas} être un contre-exemple. Il faudrait que $X_1 \to X'$ ne soit \textit{pas} une $\D$-équivalence.

\msp{page 79}

\textbf{Proposition 2.} \textit{Soit $f : X \to Y$ dans} \Dia. \textit{Conditions équivalentes :}

\hspace*{\parindent}a) \textit{Pour tout $Y'' \overset{p}{\to} Y' \to Y$, d'où}
\[
\begin{tikzcd}
X'' \arrow[r, "q"] \arrow[d] \arrow[dr, phantom, "\text{cart}"] & X' \arrow[r] \arrow[d] \arrow[dr, phantom, "\text{cart}"] & X \arrow[d, "f"] \\
Y'' \arrow[r, "p"'] & Y' \arrow[r] & Y
\end{tikzcd}
\quad ,
\]
\textit{si $p$ est une $\D$-équivalence,} \texttt{[alors]} \textit{$q$ aussi.}\footnote[35]{Il faut supposer $Y', Y''$ \textsf{Hot}-asphériques, cf. page 96, et supposer vérifié l'axiome Der 3 $a'$) de la page 97.}

\hspace*{\parindent}b) \textit{$f$ est une $\D$-fibration cohomologique ``universelle'',} i.e. \textit{est une fibration cohomologique, et le reste après tout changement de base.}

\hspace*{\parindent}b$'$) \textit{Condition duale de} (b) \textit{: $f$ est une $\D$-fibration homologique universelle.}\footnote[36]{b) $\Rightarrow$ a) \underline{canulée}, \textit{cf.} annotation marginale,\footnotemark[37] \textit{cf.} contre-exemple page 89. Mais ça marche pour les dérivateurs $\D_i$ ($i \in \mathbb{N}$) associés aux ``localisateurs fondamentaux'' $W_i$, \textit{i.e.} à la théorie des $i$-types d'homotopie.}\footnotetext[37]{L'implication b) $\Rightarrow$ a) pas prouvée, seulement sous forme plus faible :\par \hspace*{\parindent}a$'$) si $Y''$, $Y'$ sont $\D$-asphériques et 1-connexes, alors $X'' \to X'$ est une $\D$-équivalence. (Il suffit même de le supposer pour $Y'$, $Y''$ ayant des objets finaux (pour impliquer b) ou ayant des objets initiaux (pour impliquer b$'$).)}

\textit{a}) $\Rightarrow$ \textit{b}) : Supposons \textit{a}), soit $Y' \to Y$ changement de base, d'où
\[
\begin{tikzcd}
X' \arrow[r] \arrow[d, "f'"'] & X \arrow[d, "f"] \\
Y' \arrow[r] & Y
\end{tikzcd}
\quad .
\]
% foot: Version 8 juillet 2001

% ---- from Der-ChI-234_p051-055.tex
\printedpage{49}
% head: A. Grothendieck : Dérivateurs, Ch. I, Édition des Universités de Montpellier II et Paris VII
% foot: Version 8 juillet 2001

Pour prouver que $f'$ est une $\DD$-fibration cohomologique, on regarde au dessus de $Y'$ les foncteurs de la forme $\{y'\} \to Y'/y'$ ($y' \in \Ob\, Y'$) et $Y'/y'_1 \to Y'/y'_0$ (si $u' : y'_1 \to y'_0$ flèche de $Y'$), et on voit que les flèches dans $\Cat/Y'$ déduites de ceux-ci par le changement de base $X' \xrightarrow{f'} Y'$ restent des $\DD$-équivalences

\msp{page 80}

\[
\begin{tikzcd}[column sep=large, row sep=large]
X'_1 \arrow[r, "q"] \arrow[d, "f'_1"'] \arrow[dr, phantom, "\text{cart}"] & X'_0 \arrow[r] \arrow[d, "f'_0"] \arrow[dr, phantom, "\text{cart}"] & X' \arrow[r] \arrow[d, "f'"] \arrow[dr, phantom, "\text{cart}"] & X \arrow[d, "f"] \\
Y'_1 \arrow[r, "p"'] & Y'_0 \arrow[r] & Y' \arrow[r] & Y & .
\end{tikzcd}
\]
Comme $p$ est une $\DD$-équivalence, cela résulte bien de a).

\indent\textit{b}) $\Rightarrow$ \textit{a}) : Comme la condition \textit{b}) est évidemment stable par changement de base, on est ramené à prouver que si $f$ satisfait \textit{b}), et si on a
\[
\begin{tikzcd}[column sep=large, row sep=large]
X \arrow[d, "f"'] \arrow[dr, phantom, "\text{cart}"] & X' \arrow[l, "q"'] \arrow[d, "f'"] \\
Y & Y' \arrow[l, "p"]
\end{tikzcd}
\tag*{$(*)$}
\]
si $p$ est une $\DD$-équivalence, alors $q$ aussi, \textit{i.e.}
\[
\begin{tikzcd}[row sep=small]
H^*_{\DD}(X,M) \arrow[r, "\sim"] & H^*_{\DD}(X',M) \\
\wr| & \wr| \\
H^*_{\DD}(Y, f_*(M_X)) & H^*_{\DD}(Y', f'_*(M_{X'}))
\end{tikzcd}
\]
pour tout $M$ dans $\mathcal{A} = \DD(e)$. Mais $F = f_*(M_X)$ est localement constant et se calcule “fibre par fibre”, de même pour $F' = f'_*(M_{X'})$.

\msp{page 81}

L'équivalence \textit{a}) $\Leftrightarrow$ \textit{b}$'$) se déduit de \textit{a}) $\Leftrightarrow$ \textit{b}) par dualité (la condition \textit{a}) étant autoduale), ou en reprenant le même raisonnement pour \textit{a}) $\Rightarrow$ \textit{b}$'$), et pour \textit{b}$'$) $\Rightarrow$ \textit{a}) aussi (en remplaçant la cohomologie par l'homologie).

\textbf{Définition 2.} Si les conditions équivalentes de la proposition 2 sont satisfaites, on dit que $f$ est une $\DD$-\textit{fibration}.

N.B. C'est une notion \textit{à la fois} locale et “colocale” sur $Y$ (alors que fibration gauche est locale, fibration droite colocale).

\printedpage{50}
% head: A. Grothendieck : Dérivateurs, Ch. I, Édition des Universités de Montpellier II et Paris VII
% foot: Version 8 juillet 2001

\textbf{Proposition 3.} a) \textit{La notion de $\DD$-fibration est stable par changement de base et par composition.}\footnote[38]{La stabilité par composition \textit{n'est} pas claire.}

\indent b) \textit{Soient}
\[
\begin{tikzcd}[column sep=normal, row sep=large]
X \arrow[rr, "f"] \arrow[dr, "p"'] & \phantom{X} & X' \arrow[dl, "p'"] \\
& Y &
\arrow[from=1-2, to=2-2, phantom, "\text{\footnotesize comm}"]
\end{tikzcd}
\]
\textit{avec $p,p'$ des $\DD$-fibrations (il suffit que $p,p'$ soient des $\DD$-fibrations cohomologiques, ou que ce soient des $\DD$-fibrations homologiques, ou que $p,p'$ soient propres, ou $p,p'$ lisses). Si $f$ induit des $\DD$-équivalences dans les fibres, alors $f$ est une $\DD$-équivalence.}\footnote[39]{Peut se préciser considérablement dans le cas où $p$, $p'$ sont des $\DD$-fibrations, \textit{cf.} prop. 6, page 84.}

\indent\textsc{Démonstration.} \textit{a}) est immédiat sous la condition \textit{a}) (prop. 2) de la notion de $\DD$-fibration. Pour \textit{b}), supposons que $f$ induise des $\DD$-équivalences sur les fibres. On veux prouver que % [sic] On veux -> On veut, verb agreement
\[
H^*_{\DD}(X,M) \longrightarrow H^*_{\DD}(X',M)
\]
est un isomorphisme, pour tout $M$ dans $\mathcal{A}$, mais cette flèche s'identifie à
\[
H^*_{\DD}(Y, p_*(M_X)) \longrightarrow H^*_{\DD}(Y, p'_*(M_{X'}))
\]
déduite de
\[
p_*(M_X) \longrightarrow p'_*(M_{X'}) \qquad\qquad (\text{NB}\ \ M_X \simeq f^*M_{X'}\,)
\]
\msp{page 82}
et il suffit de prouver que cette flèche est un isomorphisme. Pour ceci, on regarde fibre par fibre, et on utilise le fait que “les images directes se calculent fibre par fibre”.

\textbf{Question.} Soit $f : X \to Y$ une $\DD$-fibration, et soit $F \in \DD(Y)$. Alors $f_*f^*F$ se calcule-t-il par passage aux fibres? Même question pour $f_!f^*F$. Pour la première propriété, concernant $f_*$, on montre que cela marche si $F$ est de la forme $p_*(M_Y)$, où $p : Y' \to Y$ est $\DD$-\textit{propre,} et $M \in \Ob\, \mathcal{A}$. D'autre part, la sous-catégorie pleine $C$ de $\DD(Y)$ formée des $F$ pour lesquels ça marche est (sauf erreur) “stable par $\DD$-$\mathsf{hot}$-$\varprojlim$” en le sens suivant : si on a une catégorie $I$ avec objet final $e$ formant une sous-catégorie fermée (\textit{i.e.} $e \to x \Rightarrow x = e$), et si on a $\mathcal{F} \in \DD(Y \times I)$, d'où $I^\circ \to \DD(Y)$, $i \mapsto \mathcal{F}_i$ (diagramme squelettique), si les $\mathcal{F}_i$ pour $i \neq e$ sont dans $C$, et si $\mathcal{F}$ est dans l'image essentielle de $j_*$, où $j : Y \times (I \setminus \{e\}) \hookrightarrow Y \times I$ est l'inclusion, alors $\mathcal{F}_e$ est aussi dans $C$.

\indent On peut cependant considérer que si on veux que $f : X \to Y$ satisfasse à toutes les propriétés cohomologiques et homologiques qui s'attachent à la notion d'une “fibration” (“localement triviale”), on devrait mettre dans les conditions que $f$ est \textit{lisse}, et \textit{propre}, pour que les changements de base par $f$ commutent aux foncteurs $g_*$, $g_!$. % [sic] on veux -> on veut, verb agreement

\printedpage{51}
% head: A. Grothendieck : Dérivateurs, Ch. I, Édition des Universités de Montpellier II et Paris VII
% foot: Version 8 juillet 2001

\msp{page 83}

Signalons :

\textbf{Proposition 4.} \textit{Soit $f : X \to Y$} propre. \textit{Pour que $f$ soit une $\DD$-fibration, il faut et il suffit que pour tout $y_1 \to y_0$ dans $Y$, le foncteur $X/y_1 \to X/y_0$ induit soit une $\DD$-équivalence} (i.e. \textit{que $f$ soit une $\DD$-fibration} cohomologique). \textit{Dualement, si $f$ est $\DD$-lisse, pour que ce soit une $\DD$-fibration, il faut et il suffit que pour tout $y_1 \to y_0$ dans $Y$, $y_0\backslash X \to y_1\backslash X$ soit une $\DD$-équivalence.}

\indent On pourrait appeler $\DD$-fibrations \textit{parfaites} les foncteurs à la fois propres et lisses qui sont des $\DD$-fibrations cohomologiques (ou ce qui revient au même des $\DD$-fibrations homologiques).

\textbf{Corollaire.} \textit{Supposons $f$ $\Cat$-cofibration (donc propre). Pour que $f$ soit une $\DD$-fibration, il faut et il suffit que les foncteurs $X_{y_0} \to X_{y_1}$ (pour $y_0 \to u_1$ dans $Y$) soient des $\DD$-équivalences. Duallement, pour $\Cat$-fibration (donc lisse).} % [sic] u_1 -> y_1, the arrow is y_0 -> y_1 ; [sic] Duallement -> Dualement, spelling

\textbf{Proposition 5.} \textit{Soit $f : X \to Y$ dans} $\Dia$. \textit{Conditions équivalentes :}
\begin{enumerate}[label={}, leftmargin=2em, labelsep=0pt]
\item a) \textit{$f$ est une $\DD$-équivalence, et le reste par tout changement de base ;}
\item b) \textit{$f$ est une $\DD$-fibration à fibres $\DD$-contractiles.}
\end{enumerate}

\textsc{Démonstration.} \textit{a}) $\Rightarrow$ \textit{b}) : Que \textit{a}) implique que $f$ soit une $\DD$-fibration est immédiat sur la définition de celle-ci (condition \textit{a}) de la proposition 2), et que les fibres sont contractiles est tautologique.

\textit{b}) $\Rightarrow$ \textit{a}) : Il suffit de montrer que \textit{b}) implique que $f$ est une $\DD$-équivalence. Mais c'est un cas particulier de la proposition 3 \textit{b}) (en y prenant $X' = Y$).

\msp{page 84}

Quand ces conditions sont satisfaites, on s'attend à ce que le foncteur
\[
f^* : \DD(Y) \longrightarrow \DD(X)
\]
soit pleinement fidèle, \textit{i.e.} que les morphisme d'adjonction % [sic] les morphisme -> les morphismes, plural agreement
\[
F \longrightarrow f_*f^*F \quad , \qquad f_!f^* \longrightarrow F % [sic] f_!f^* -> f_!f^*F, the object F is missing
\]
dans $\DD(Y)$ soient des isomorphismes. (C'est vrai si $f$ est $\DD$-propre ou $\DD$-lisse, cf. cor. 3, p.71.) Cela revient encore à dire que les $f_*f^*F$ se calculent par passage aux fibres, ou encore itou pour les $f_!f^*F$. Cela nous ramène à la question de la page 82.

\textbf{Définition 3.} On dit alors que $f$ est une $\DD$-fibration \textit{triviale}.

\textbf{Proposition 6} (précise prop. 3 page 81). \textit{Considérons un triangle commutatif dans} $\Dia$
\[
\begin{tikzcd}[column sep=small, row sep=large]
X \arrow[rr, "f"] \arrow[dr, "p"'] & & X' \arrow[dl, "p'"] \\
& S & & & .
\end{tikzcd}
\]

\printedpage{52}
% head: A. Grothendieck : Dérivateurs, Ch. I, Édition des Universités de Montpellier II et Paris VII
% foot: Version 8 juillet 2001

\indent a) \textit{Supposons que $p$ et $p'$ soient des $\DD$-fibrations gauches, ou que $p$ et $p'$ soient des $\DD$-fibrations droites. Soit $S_0$ une partie de $S$ qui rencontre toute composante connexe de $S$. Pour que $f$ soit une $\DD$-équivalence, il} suffit \textit{que pour tout $s_0 \in S_0$, $f$ induise une $\DD$-équivalence $X_{s_0} \to Y_{s_0}$.}

\indent b) \textit{Cette condition est aussi nécessaire}\footnote[40]{Pas prouvé!\footnotemark[41] Même si $X' = S$, ce n'est pas clair, même si $\DD = \mathsf{HOTAB}$, $f$ propre et lisse. Mais OK si $S$ est $\DD$-contractile.}\footnotetext[41]{\underline{Faux}, \textit{cf.} contre-exemple, p. 87.} \textit{si $p$ et $p'$ sont des $\DD$-fibrations. Donc dans ce cas, $f$ est une $\DD$-équivalence si et seulement si $f$ induit des $\DD$-équivalences fibre par fibre.}

\msp{page 85}

\textsc{Démonstration.} a) Il suffit, en vertu de la proposition 3, de montrer que l'hypothèse implique que pour \textit{tout} $s \in S$, $f$ induit une $\DD$-équivalence $f_s : X_s \to X'_s$. Cela résulte aussitôt du lemme :

\textbf{Lemme.} \textit{Si $p$ et $p'$ sont toutes deux soit des $\DD$-fibrations gauches, soit des $\DD$-fibrations droites, et si $s_1 \to s_0$ est une flèche dans $S$, alors $f_{s_1}$ est une $\DD$-équivalence si et seulement si $f_{s_0}$ l'est. (Donc l'ensemble $S' \subset \Ob\, S$ des $s \in S$ tels que $f_s$ soit une $\DD$-équivalence est une réunion de composantes connexes de $S$.)}

\indent Supposons par exemple que $p,p'$ soient des $\DD$-fibrations gauches. Considérons, pour $s \in S$, le diagramme
\[
\begin{tikzcd}[row sep=large]
X_s \arrow[r, "f_s"] \arrow[d] & X'_s \arrow[d] \\
X/s \arrow[r, "f/s"'] & X'/s & .
\end{tikzcd}
\]
Les flèches verticales sont des $\DD$-équivalences par l'hypothèse sur $p,p'$, donc $f_s$ est une $\DD$-équivalence si et seulement si $f/s$ l'est. Considérons maintenant
\[
\begin{tikzcd}[row sep=large]
X/s_1 \arrow[r, "f/s_1"] \arrow[d] & X'/s_1 \arrow[d] \\
X/s_0 \arrow[r, "f/s_0"'] & X'/s_0 & .
\end{tikzcd}
\]
Ici encore, l'hypothèse sur $p,p'$ implique que les flèches verticales sont des $\DD$-équivalences, donc $f_{s_1}$ est une $\DD$-équivalence si et seulement si $f_{s_0}$ l'est.

\msp{page 86}

Si $Y' = S$, l'énoncé (non prouvé) dit que si $f : X \to S$ est une $\DD$-fibration et une $\DD$-équivalence, alors les fibres sont $\DD$-contractiles. Ce n'est pas sûr même pour une projection $X = T \times S \to S$. Si $\eta = \mathit{hot}_0(S)$, $\xi = \mathit{hot}_0(T)$, la question se réduit à celle-ci : Si pour tout $M$ dans $\mathcal{A}$, on a : l'homomorphisme
\[
\alpha_M : \xi \times_{\DD} M \longrightarrow M
\]
\printedpage{53}
% head: A. Grothendieck : Dérivateurs, Ch. I, Édition des Universités de Montpellier II et Paris VII
% foot: Version 8 juillet 2001
satisfait
\[
\eta \times_{\DD} \alpha_M : \eta \times_{\DD} (\xi \times_{\DD} M) \xrightarrow{\sim} \eta \times_{\DD} M
\]
est un isomorphisme, alors $\alpha_M$ est lui même un isomorphisme.

\msp{page 87}

Voici un \textit{contre-exemple}\footnote[42]{Voir un contre-exemple encore plus simple pages 93, 94. (Mais il a l'inconvenient de ne pas être à fibres connexes.) % [sic] inconvenient -> inconvénient, missing accent
} que je donne d'abord dans le contexte des espaces topologiques, au lieu de $\Cat$. Je vais construire une fibration localement triviale et non triviale
\[
\begin{tikzcd}[column sep=0.2em, row sep=large]
X \arrow[d, "f"'] & \supset S^1 & \\
S & = S^1 & \qquad ,
\end{tikzcd}
\]
de fibre $S^1$, telle que $f$ soit une $\DD$-équivalence pour tout $\DD = \mathsf{Der}(k\text{-}\mathsf{Mod})$, $k$ étant un anneau où 2 est inversible. Cela signifie que
\[
H_*(X,k) \longrightarrow H_*(S,k)
\]
est un isomorphisme modulo 2-torsion (noyau et conoyau sont annulés par une puissance de 2 - en fait par 2 ici), ou encore $H_*(X,k) \xrightarrow{\sim} H_*(S,k)$ si 2 inversible dans $k$. Mais la fibre $S^1$ n'est pas $\DD$-contractile (sauf si $k = 0$!).\footnote[43]{Exemple beaucoup plus simple encore : on prend $X = S^1$, et $X \to S$ le revêtement à deux feuillets. Si on prend le revêtement à $n$ feuillets, c'est aussi OK si $k$ anneau où $n$ est inversible.}

\indent La fibration est donnée à isomorphisme unique près, en se fixant un point $s \in S$, et en se donnant un automorphisme $u$ de la fibre $X_s$ (“monodromie”). (La classe d'isomorphie du fibré ne dépend que de la classe d'isotopie de $u$). Je prend pour $u$ la symétrie (inversant l'orientation) $z \mapsto z^{-1}$ du cercle standard $\{z \in \mathbb{C} \mid |z| = 1\}$. Cette fibration “se réalise” par des ensembles ordonés $X \to S$, $S^1$ étant représenté (sur la base, comme sur les fibres) par % [sic] Je prend -> Je prends, verb agreement ; [sic] ordonés -> ordonnés, spelling
\[
\begin{tikzcd}[column sep=normal, row sep=normal]
& & a & & \\
S \ = & 0 \arrow[ur] \arrow[dr] & & 0' \arrow[ul] \arrow[dl] & = \ S' \amalg_R S'' \quad , \\
& & b & &
\end{tikzcd}
\]
où $R$ est la catégorie discrète ayant deux éléments $a$, $b$. On définit $X$ comme $X' \amalg_Y X''$, où $X' = X|S'$, $X'' = X|S''$, $Y = X|R$.\footnote[44]{N.B. $S'$, $S''$, $R = S' \cap S''$ fermés dans $S$, $X'$, $X''$, $Y = X' \cap X''$ fermés dans $X$.}

% ---- from Der-ChI-234_p056-060.tex
\printedpage{54}
% head: A. Grothendieck : Dérivateurs, Ch. I, Édition des Universités de Montpellier II et Paris VII

\msp{page 88}
On pose
\[ X' = S' \times \bar S \qquad\qquad X'' = S'' \times \bar S \]
($\bar S$ copie conforme de $S$, cf ci-dessous) et il faut, pour recoller $X'$ et $X''$, se donner un isomorphisme entre $X'|R \simeq R \times \bar S = \bar S \amalg \bar S$ et $X''|R \simeq R \times \bar S = \bar S \amalg \bar S$ qui soit un $R$-foncteur, \textit{i.e.} qui provient de deux isomorphismes
\[
\begin{aligned}
\bar S &\xrightarrow{u_a} \bar S \\
\bar S &\xrightarrow{u_b} \bar S \qquad .
\end{aligned}
\]
On prend
\[
\begin{aligned}
u_a &= \mathsf{id}_{\bar S} \\
u_b &= \text{symétrie par rapport à l'axe } \bar 0, \bar 0' \ \textit{i.e.} % Doubt resolved: PDF leaf 18 (printed p.54) reads u_b, with plain b.
\\
&\phantom{=}\quad u_b \text{ laisse fixes } \bar 0, \bar 0' \text{ et intervertit } \bar a \text{ et } \bar b.
\end{aligned}
\]
(C'est plus clair si on désidentifie la base et la fibre, en prenant fibre
\[
\bar S \quad = \quad
\begin{tikzcd}
 & \bar a & \\
\bar 0 \arrow[ur] \arrow[dr] & & \bar 0' \arrow[ul] \arrow[dl] \\
 & \bar b &
\end{tikzcd}
\qquad ).
\]

Pour voir que $X$ est une $\D$-fibration sur $S$ (quelque soit $\D$) prouvons qu'il est \textit{propre}\footnote[45]{N.B. $f : X \to S$ est également \textit{lisse}, car c'est un fibré trivial (un produit) au dessus de deux ouverts $S - \{a\}$, $S - \{b\}$ qui recouvrent $S$. Mieux : $f$ est à la fois $\Cat$-fibrant et $\Cat$-cofibrant (car cette condition n'utilise que les propriétés des $X \times_S \boldsymbol{\Delta}_2 \to \boldsymbol{\Delta}_2$, pour $\boldsymbol{\Delta}_2 \to S$).} sur $S$ - on est ramenés à voir que $X' \to S'$ et $X'' \to S''$ sont propres, ce qui est trivial. Pour voir que c'est une $\D$-fibration, on est ramenés également au cas de $X' \to S'$ et $X'' \to S''$, où c'est aussi trivial. On a utilisé le : % [sic] quelque soit -> quel que soit, spelling

\textbf{Lemme.} \textit{Soient $f : X \to S$, $(S_i)_{i \in I}$ une famille de sous-catégories fermées recouvrant $S$. Pour tout $i$, soit $X_i = X \times_S S_i \xrightarrow{f_i} S_i$. Pour que $f$ soit propre, il faut et il suffit que les $f_i$ le soient. Pour que $f$ soit une $\D$-fibration propre il faut et il suffit que les $f_i$ le soient.}

\msp{page 89}
Voici le petit calcul qui montre que $f : X \to S$ est une $\D$-équivalence. On fait le calcul en cohomologie.\footnote[46]{Mais dans le cas du revêtement à 2 feuillets de $S$, il n'y a pas même de calcul à faire!} On doit prouver que
\[ H^*(S,k) \longrightarrow H^*(X,k) \tag*{$(*)$} \]
\printedpage{55}
% head: A. Grothendieck : Dérivateurs, Ch. I, Édition des Universités de Montpellier II et Paris VII
est un isomorphisme (et de même pour tout coefficient constant $M$ sur $S$). Or
\[ H^*(X,k) \overset{\mathrm{def}}{=} \RR\Gamma_X(k_X) = \RR\Gamma_S(\RR f_*(k_X)) \quad , \]
et $(*)$ est induit par
\[ i : k_S \longrightarrow \RR f_* k_X \quad . \]
Or on a
\[ \RR f_* k_X \simeq \underbrace{k_S}_{= \mathsf{Im}(i)} \quad \oplus \quad \underbrace{k_S[1]}_{k_S \text{ en degré } 1}{}^{\underline{\mathrm{tordu}}} \]
savoir ``tordu'' par l'action de monodromie sur $S = S^1$, laquelle est la symétrie $\lambda \mapsto -\lambda$. On a
\[ \RR\Gamma_S(k_S[1]^{\mathrm{tordu}}) = H^1(\underbrace{\pi_1(S)}_{= \mathbb{Z}}\, , k_S^{\mathrm{tordu}}) = k_{(\mathbb{Z})} = k/(u - \mathsf{id})k = k/2k = 0 \ , \]
d'où le résultat.

Notons que si on regarde
\[ X' = X \times_S X \xrightarrow{f'} S \quad , \]
on a par Künneth relatif
\[ \RR f'_*(k_{X'}) \simeq \RR f_*(k_X) \otimes^{\mathbb{L}} \RR f_*(k_X) \simeq k_S \oplus k_S[1]^t \oplus k_S[1]^t \oplus \underbrace{k_S[1]^t \otimes k_S[1]^t}_{\simeq k_S[2] \text{ non tordu}} \ , \]
d'où
\[ H^*_{\D}(X',k) \simeq \RR\Gamma_{X'}(k_{X'}) \simeq \RR\Gamma_S(\RR f'_*(k_{X'})) \simeq \RR\Gamma_S(k_S) \oplus \RR\Gamma_S(k_S)[2] \]
et
\[ H^*_{\D}(S,k) \longrightarrow H^*_{\D}(X',k) \]
n'est \textit{pas} un isomorphisme.

\msp{page 90}
Ainsi $f : X \to S$ qui est une $\D$ fibration (propre et lisse) satisfait bien les conditions équivalentes $b$), $b'$) de la proposition 2 (page 79), mais pas la condition $a$), puisque le changement de base
\[
\begin{tikzcd}[column sep=large, row sep=large]
X' = X \times_S X \arrow[r, "{q = pr_1}"] \arrow[d, "{f' = pr_2}"'] & X \arrow[d, "f"] \\
S' = X \arrow[r, "{p = f}"'] & S
\end{tikzcd}
\]
est une $\D$-équivalence $S' \to S$, alors que $X' \to X$ ne l'est pas.

Voici un autre contre-exemple, sur le même principe mais un peu plus simple. Soit $G$ un groupe fini (par exemple $G = \mathbb{Z}/2\mathbb{Z}$), faisons le opérer sur un objet $Z$ de $\Cat$, soit $S = [G] = B_G$ la catégorie réduite à un objet $a$, avec $G = \mathsf{End}(a)$. Soit $E_G$ la catégorie (isomorphe à $(B_G)/a$)
\[ \mathsf{Ob}\ E_G = G \quad , \qquad \mathsf{Hom}_{E_G}(g,g') = \{ h \in G \mid g = g'h \} \quad , \]
composition évidente - alors $E_G$ est un groupoïde trivial (1-connexe). Le goupe $G$ y opère à gauche librement par translations, et $B_G = E_G/G$. % [sic] goupe -> groupe, misprint

\printedpage{56}
% head: A. Grothendieck : Dérivateurs, Ch. I, Édition des Universités de Montpellier II et Paris VII

\msp{page 91}
Si $G$ opère à gauche sur $Z$, on peut alors l'utiliser pour tordre $Z$, par
\[ X = E_G \times_G Z = (E_G \times Z)/\text{opérations ``diagonales'' de G} \quad . \]
Donc on a un diagramme cartésien
\[
\begin{tikzcd}
 & X' \arrow[d, equal] & X'' \arrow[d, equal] \\
X \arrow[d, "f"'] & Z \times E_G \arrow[l] \arrow[d, "pr_2"] & Z \arrow[l, dashed, "{\overset{\D}{\sim}}"'] \arrow[d] \\
B_G \arrow[d, equal] & E_G \arrow[l] \arrow[d, equal] & e \arrow[l, dashed, "{\overset{\D}{\sim}}"'] \arrow[d, equal] \\
S & S' & S''
\end{tikzcd}
\]
Comme le changement de base $S' = E_G = (B_G)/a \to B_G = S$ transforme $X$ sur $B_G$ en un produit (lequel est un $\D$-fibré propre et lisse sur $S'$), on voit que $X$ sur $B_G$ est un $\D$-fibré propre et lisse (car c'est là une condition locale sur $B_G$).

On suppose maintenant que $G$ est fini d'ordre $n > 1$, et soit $k$ un anneau où $n$ est invesible, $\D = \mathsf{Der}(k - \mathsf{Mod})$. Alors $E_G \to B_G$ est une $\D$-équivalence, ou ce qui revient au même $e \to B_G$ l'est.\footnote[47]{En fait, $e$, $E_G$, $B_G$ sont $\D$-asphériques.} Pourtant, après changement de base soit par $S' = E_G$, soit par $S'' = e$, $X' \to X$ (ou ce qui revient au même, $X'' \to X$) n'a guère de raison d'en être aussi. On a en cohomologie % [sic] invesible -> inversible, misprint
\[ \RR\Gamma_X(k_X) \simeq \RR\Gamma_{BG}(f_*(k_X)) \simeq \RR\Gamma_G(G, H^*(Z,k)) \]
\msp{page 92}
donc ($n$ étant inversible dans $k$) :
\[ H^*(X,k) = \RR\Gamma_X(k_X) \simeq H^*(Z,k)^G \longrightarrow H^*(Z,k) \quad , \]
ce n'est pas en général un isomorphisme. On peut prendre par exemple $G = \mathbb{Z}/2\mathbb{Z}$, $Z \sim S^1$ comme avant, et $G$ opérant par symétrie de $S^1$ comme précédemment.

Cet exemple est particulièrement convainquant, car c'est un exemple \textit{avec $S$, $S'$ $\D$-asphériques,} et $S' \to S$ un isomorphisme local (en fait, un revêtement principal de $S$). C'est aussi un exemple d'un $\D$-fibré lisse et propre $X \to S$, sur une base \textit{asphérique} $S$, et tel que l'inclusion d'une fibre dans $X$ % [sic] convainquant -> convaincant, spelling of the adjective
\[ X_s \longrightarrow X \]
ne soit \textit{pas} une $\D$-équivalence.

Le seul reproche qu'on puisse faire à cet exemple, en comparaison avec le précédent, c'est que les catégories $X, S$ ne sont pas des ensembles ordonnés. Et $S$ ne peut être un ensemble ordonné fini, car l'espace classifiant d'un groupe fini a de la cohomologie en toutes dimensions (sauf erreur).

\printedpage{57}
% head: A. Grothendieck : Dérivateurs, Ch. I, Édition des Universités de Montpellier II et Paris VII

\msp{page 93}
Mais on peut s'inspirer de cet exemple pour le modifier de façon qu'il se transpose à des ensembles ordonnés. Il suffit de réaliser, pour $G$ donné :

\textit{a}) Une opération \textit{libre} de $G$ sur un objet $E$ de $\Cat$, telle que $H^*(E,k)^G \simeq k$, \textit{i.e.} $B = E/G$ soit $\D$-asphérique.

\textit{b}) Une opération de $G$ sur un objet $Z$, telle que l'inclusion
\[ H^*(Z,k)^G \longrightarrow H^*(Z,k)^G \] % [sic] H^*(Z,k)^G \longrightarrow H^*(Z,k)^G -> H^*(Z,k)^G \longrightarrow H^*(Z,k), an inclusion of the invariants (cf. p. 56)
ne soit pas \texttt{[un]} isomorphisme.

\textit{c}) Choisir pour $E$, $Z$ des ensembles ordonnés, et s'assurer que $B = E/G$ est également représenté par une catégorie ordonnée. Sauf erreur, il en sera de même pour $X = E \times Z/G$. On aura en tout les cas % [sic] en tout les cas -> en tous les cas, agreement
\[
\begin{tikzcd}[column sep=large]
 & Z \times E \arrow[d, equal] & \\
X \arrow[d] & X' \arrow[l] \arrow[d] & Z \arrow[l] \arrow[d] \\
S \arrow[d, equal] & S' \arrow[l] \arrow[d, equal] & S'' = \{s\}\, , \quad s \in E \ . \arrow[l] \\
E/G & E & \text{\footnotesize asph} \\
\text{\footnotesize asph} & \substack{\text{asph} \\ \text{si on peut}} &
\end{tikzcd}
\]
Cela donne encore un exemple avec $X$ $\D$-fibré propre et lisse, sur une base $S$ $\D$-asphérique, tel que l'inclusion d'une fibre $X_s$ dans $X$ ne soit pas une $\D$-équivalence. On peut même se débrouiller pour que $H^*(Z,k)^G \simeq k$ ce qui implique
\msp{page 94}
que $X$ est lui aussi asphérique (mais les fibres de $X \to S$ ne sont pas $\D$-asphériques). Cela donne donc un autre contre-exemple au b) de la prop. 6 (p. 84).

Voici la façon la plus économique, je crois, de réaliser la situation que j'ai dite. Topologiquement, je prend % [sic] je prend -> je prends, conjugation
\[
\begin{aligned}
G &= \{\pm 1\} \\
E &= S^2 &\qquad& \text{sphère euclidienne, sur laquelle } \{\pm 1\} \\
  &      &&      \text{opère par antipodisme } x \mapsto -x \\
B &= S^2/\{\pm 1\} && \text{plan projectif réel} \\
Z &= S^1 && \text{où G opère comme précédemment par la symétrie\,.}
\end{aligned}
\]
C'est donc un exemple très proche de celui du début (p. 87), où on pouvait aussi considérer la base ($B = S^1$) comme quotient d'un $E$ ($= S^1$) par $G = \{\pm 1\}$ opérant librement (par
\printedpage{58}
% head: A. Grothendieck : Dérivateurs, Ch. I, Édition des Universités de Montpellier II et Paris VII
l'antipodisme $z \mapsto -z$ dans $\mathbb{C} \simeq \RR^2$). Mais dans cet exemple, la base n'était pas $\D$-asphérique.

La façon la plus économique de réaliser la sphère $S^2$ par un ensemble ordonné est sans doute par sa décomposition ``globulaire'' par trois paires de cellules de dimension 0, 1, 2, donnant l'ensemble ordonné
\[
I \qquad = \qquad
\begin{tikzcd}[row sep=large, column sep=large]
S_2^- & S_2^+ \\
S_1^- \arrow[u] \arrow[ur] & S_1^+ \arrow[u] \arrow[ul] \\
S_0^- \arrow[u] \arrow[ur] & S_0^+ \arrow[u] \arrow[ul]
\end{tikzcd}
\]
et l'antipodisme échange $S_i^-$ et $S_i^+$, mais le quotient n'est pas ordonné, c'est la catégorie (importante pour la théorie des 2-catégories) :
\msp{page 95}
\[
\begin{tikzcd}[row sep=large]
S_2 \\
S_1 \arrow[u, shift left, "p_1^-"] \arrow[u, shift right, "p_1^+"'] \\
S_0 \arrow[u, shift left, "p_0^-"] \arrow[u, shift right, "p_0^+"']
\end{tikzcd}
\qquad\qquad \text{avec} \qquad\qquad
\begin{cases}
p_1^- p_0^- = p_1^- p_0^+ \\
p_1^+ p_0^- = p_1^+ p_0^+
\end{cases}
\]
Mais prenons la subdivision barycentrique de la décomposition cellulaire précédente de $S$ dont les sommets correspondent aux éléments de $\mathsf{Ob}\ I$, les simplexes aux parties totalement ordonnés de $\mathsf{Ob}\ I$. On trouve l'octaèdre, ayant 6 sommets, 12 arêtes, 8 faces, et $G$ opère librement sur l'ensemble des simplexes, donc le quotient : 3 sommets, 6 arêtes, 4 faces, décomposition cellulaire du plan projectif. (Sa décomposition par trois droites.) On présente $S^1$ par l'ensemble ordonné (comme au début) % [sic] parties totalement ordonnés -> parties totalement ordonnées, agreement
\[
\begin{tikzcd}[row sep=large, column sep=large]
S_1^- & S_1^+ \\
S_0^- \arrow[u] \arrow[ur] & S_0^+ \arrow[u] \arrow[ul]
\end{tikzcd}
\qquad\qquad \textit{i.e.} \qquad\qquad
\begin{tikzcd}
 & S_1^+ & \\
0 = S_0^- \arrow[ur] \arrow[dr] & & S_0^+ = 0' \arrow[ul] \arrow[dl] \\
 & S_1^- &
\end{tikzcd}
\]
et $G$ y opère par symétrie échangeant $S_1^-$ et $S_1^+$. Donc l'ensemble ordonné $X$ aura $4 \times 13 = 52$ sommets et c'est modeste \ldots

% ---- from Der-ChI-234_p061-065.tex
\printedpage{59}
% head: A. Grothendieck : Dérivateurs, Ch. I, Édition des Universités de Montpellier II et Paris VII
% foot: Version 8 juillet 2001

\msp{page 96}
Finalement, il semble qu’il ne reste plus rien de la proposition 2 (p.79) introduisant la notion de $\DD$-fibration – pas même l’équivalence des deux conditions $b$) et $b'$) ($\DD$-fibration à gauche universelle, $\DD$-fibration à droite universelle). On songe à leur relation à une condition $a$) modifiée ainsi (pour $f : X \to Y$) :

$a$) Pour tout $Y'' \to Y' \to Y$, avec $Y'$ et $Y''$ \Hot-\underline{asphériques} (et pas seulement $\DD$-asphériques) $X'' = X \times_Y Y'' \to X' = X \times_Y Y'$ est une $\DD$-équivalence. (Quand on suppose que $Y''$ et $Y'$ ont des éléments finaux, c’est la condition que $f$ soit une $\DD$-fibration à gauche universelle; quand on suppose que $Y''$ et $Y'$ ont des éléments initiaux, c’est la condition $\DD$-fibration à droite universelle.)\footnote[48]{Il faut expliciter le raisonnement. Il suffit de tester même sur $Y'' \to Y'$, où $Y''$ est \underline{ponctuelle}, et où $Y'$ a un objet final (resp. initial).}

On a évidemment $a$) $\Rightarrow$ $b$) et $a$) $\Rightarrow$ $b'$) et on aimerait que ce soient des équivalences. Voici un axiome intéressant\footnote[49]{\textit{cf.} \textbf{Der 3} $a'$) page suivante.} sur $\DD$, qui semble vérifié dans tous les cas connus, qui assure l’équivalence souhaitée. Je vais en même temps reformuler \textbf{Der 3} en trois volets (\textit{cf.} p. 67).

\msp{page 97}
\textbf{Der} 3 $a$) Si $Z \in \Dia$ est \Hot-asphérique, alors $p : Z \to e$ est une $\DD$-équivalence, \textit{i.e.} par exemple $p^* : \mathcal{A} = \DD(e) \to \DD(Z)$ est pleinement fidèle.

$b$) Soit $f : Z \to Z'$ une flèche dans \Dia. Si $f$ est une \Hot-équivalence, c’est une $\DD$-équivalence (\textit{i.e.} $f^* : \DD(Z') \to \DD(Z)$ induit une \underline{équivalence} entre les sous-catégories de $Z$, $Z'$ \texttt{[plutôt de} $\DD(Z)$\texttt{,} $\DD(Z')$\texttt{]} formées des objets constants).

$a'$) (Précise considérablement $a$). Si $Z \in \Dia$ est \Hot-asphérique, alors on a $a$), et l’image essentielle de $\DD(e) = \mathcal{A}$ dans $\DD(Z)$ est formée des objets \underline{localement} \underline{constants} de $Z$, \textit{i.e.} des objets tels que le foncteur
\[
Z^{\circ} \longrightarrow \mathcal{A} \qquad z \longmapsto F_z
\]
soit localement constant (\textit{i.e.} transforme toute flèche en isomorphisme) donc isomorphe à un foncteur constant (puisque $Z$ est \Hot-asphérique, donc 1-connexe).

Il faut bien noter que dans $a'$), il ne suffit \underline{pas} de supposer que $Z$ soit $\DD$-asphérique, au lieu de \Hot-asphérique. Dès que la $\DD$-équivalence est strictement plus faible que \Hot, il semble que $a'$) devient toujours faux, en tant que condition sur $\DD$, si on y suppose seulement que $Z$ soit $\DD$-asphérique. Par exemple, même si $\DD = \DD_i$ (théorie des $i$-types d’homotopie relative), la catégorie
\msp{page 98}
des objets localement constants sur $Z$ s’exprime en termes du $(i+1)$-type d’homotopie, mais pas en termes de son $i$-type d’homotopie seul.

Il semble bien que ce soit l’axiome \textbf{Der} 3 ($a'$) qui assure l’existence d’une bonne notion (“bilatère”) de $\DD$-fibration, qui soit stable par changement de base, et qui se réduise soit
\printedpage{60}
% head: A. Grothendieck : Dérivateurs, Ch. I, Édition des Universités de Montpellier II et Paris VII
% foot: Version 8 juillet 2001
à “$\DD$-fibration à gauche universelle” soit à “$\DD$-fibration à droite universelle” et donc qui soit en plus une condition \underline{vérifiable} pratiquement. Dans le cas où $f$ est propre ou lisse, le critère est particulièrement simple, puisque il se borne à la condition que les $f_* f^* a_Y$, ou les $f_! f^* a_Y$ (pour $a \in \mathcal{A}$), soient des objets localements constants sur $Y$ (donc \underline{constants} si $Y$ est \underline{\underline{\Hot}}-asphérique, en vertu de \textbf{Der} 3 ($a'$). % [sic] localements -> localement, plural ending on the adverb; % [sic] "(donc ... ($a'$)." -> closing parenthesis after "($a'$)" missing, unbalanced parenthesis

Je vais également introduire une variante renforcée de \textbf{Der} 3 $b$) :

\textbf{Der} 3 $b'$) Soit $Z' \to Z$ une \Hot-équivalence dans \Dia. Alors $f^* : \DD(Z) \to \DD(Z')$ induit une équivalence entre les catégories des coefficients \underline{localements} \underline{constants} sur $Z$ et sur $Z'$.\footnote[50]{\textbf{Der} 3 $b'_0$) : le foncteur est fidèle\\
\hspace*{6em}$b'_1$) : le foncteur est pleinement fidèle\\
\hspace*{6em}$b'_2$) : le foncteur est une équivalence} % [sic] localements -> localement, plural ending on the adverb

\msp{page 99}
Dans certaines questions il suffit de savoir que le dit foncteur est pleinement fidèle. Cela implique que
\[
H^*_{\DD}(Z,F) \overset{\sim}{\longrightarrow} H^*_{\DD}(Z', f^*F)
\]
quand $F$ est \underline{localement} \underline{constant} sur $Z$\footnote[51]{On pourra dire que $f : Z' \to Z$ est une $\DD$-équivalence 0-\textit{complète à gauche,} et 0-\textit{complète à droite} dans le cas homologique.} (mais il n’est pas clair que cette condition implique toujours la pleine fidélité\footnote[52]{Mais c’est OK si on a des \underline{\Hom} internes dans $\DD(Z)_{\mathrm{loc.const.}}$ satisfaisant
\[
\Hom(F,G) \simeq \Hom(e_Z, \underline{\Hom}(F,G)) \quad ,
\]
où $e$ est l’élément final de $\mathcal{A}$, les \underline{\Hom} étant compatibles aux images inverses.}) et de même pour l’homologie
\[
H_*^{\DD}(Z', f^*F) \longrightarrow H_*^{\DD}(Z,F) \quad .
\]
Moyennant cette condition, on a la :

\textbf{Proposition :} \textit{On suppose} \textbf{Der} 3 \textup{b}$'$) \textit{satisfait. Soit $f : X \to Y$ dans} \Dia. \textit{Pour que $f$ soit une $\DD$-fibration, il faut et il suffit que pour $Y'' \to Y'$ sur $Y$ qui soit une} \Hot\textit{-équivalence, $X'' = X \times_Y Y'' \to X' = X \times_Y Y'$ soit une $\DD$-équivalence.}

À vrai dire, on peut énoncer quelque chose de raisonnable, sans supposer l’axiome \textbf{Der} 3 $b'$). On doit introduire la :

\textbf{Définition :} Soit $f : Z' \to Z$ dans \Dia. On dit que $f$ et une $\DD$-équivalence 1-\textit{complète} si le foncteur $f^* : \DD(Z) \to \DD(Z')$ est pleinement fidèle sur la sous-catégorie pleine de $\DD(Z)$ formée des coefficients localement constants. On dit que $f$ est une $\DD$-équivalence 2-\textit{complète} si elle induit une équivalence entre les catégories des coefficients localement constants. % [sic] f et une -> f est une, verb "est" misprinted as "et"

\msp{page 100}
On dit que $Z$ est \textit{complètement} $\DD$-\textit{asphérique}, si $Z \to e$ est une $\DD$-équivalence 2-complète.

\printedpage{61}
% head: A. Grothendieck : Dérivateurs, Ch. I, Édition des Universités de Montpellier II et Paris VII
% foot: Version 8 juillet 2001

Ainsi, l’axiome \textbf{Der} 3 $a'$) peut se reformuler, en disant que si $Z$ est \Hot-asphérique, alors $Z$ est complètement $\DD$-asphérique. Et de même, \textbf{Der} 3 $b'$) en disant que si $f : X \to Y$ dans \Dia{} est une \Hot-équivalence, alors $f$ est une $\DD$-équivalence 2-complète. Nous allons donner une variante affaiblie, de \textbf{Der} 3 $a'$), qui suffit pour donner une notion raisonnable de $\DD$-fibration :

\textbf{Der} 3 $a''$) Tout $Z$ dans \Dia{} qui a soit un objet final, soit un objet initial, est complètement $\DD$-asphérique.

\textbf{Définition.} On dit que $f : X \to Y$ dans \Dia{} est une $\DD$-fibration, si pour
\[
Y'' \longrightarrow Y' \longrightarrow Y
\]
avec $Y'' \to Y'$ une $\DD$-équivalence 1-complète (par exemple $Y'$ et $Y''$ complètement $\DD$-asphèriques), $X'' = X \times_Y Y'' \longrightarrow X' = X \times_Y Y'$ est une $\DD$-équivalence. % [sic] asphèriques -> asphériques, grave accent instead of acute

Cela équivaut à chacune des conditions $b$), $b'$) de la prop. 2 (page 79), \textit{i.e.} que $f$ soit universellement une $\DD$ fibration à gauche, ou universellement une $\DD$-fibration à droite (grâce à \textbf{Der 3} $b''$ \texttt{[plutôt} $a''$\texttt{]}). % [sic] $\DD$ fibration -> $\DD$-fibration, hyphen missing

\msp{page 101}
\textbf{Définition 1.} $f : X \to Y$ dans \Dia{} est appelé $\DD$-\textit{équivalence cohomologique} (sous-entendu : \textit{complète}) si pour tout $F \in \Ob\,\DD(Y)$, $F$ \underline{localement} \underline{constant}, on a
\[
H^*_{\DD}(Y,F) \overset{\sim}{\longrightarrow} H^*_{\DD}(X, f^*(F)) \quad .
\]
Notion duale $\DD$-\textit{équivalence homologique}
\[
H_*^{\DD}(X, f^*(F)) \overset{\sim}{\longrightarrow} H_*^{\DD}(Y,F) \quad .
\]
Si $Y$ est complètement $\DD$-asphérique, alors les coefficients localement constants sur $Y$ sont constants, et $f$ est une $\DD$-équivalence cohomologique (resp. homologique) si et seulement si c’est une $\DD$-équivalence \textit{i.e.} si $X$ est $\DD$-asphérique.\footnote[53]{L’axiome \textbf{Der} 3 $b'_1$) implique que toute \Hot-équivalence est une $\DD$-équivalence cohomologique.}

\textbf{Proposition 1} (Soritale). \textit{Soient $X \overset{f}{\to} Y \overset{g}{\to} Z$. Si $f$ et $g$ sont des $\DD$-équivalences cohomologiques, alors $gf$ aussi. Si $gf$ et $f$ sont des $\DD$-équivalences cohomologiques, alors $g$ aussi. Si $gf$ est une $\DD$-équivalence cohomologique, et si $g$ est une $\DD$-équivalence \underline{2-complète}} (donc tout $F$ localement constant sur $Y$ provient d’un $G$ localement constant sur $Z$), \textit{alors $f$ est une $\DD$-équivalence cohomologique.}

\textbf{Définition 2.} 1°) $f : X \to Y$ dans \Dia{} est appelé une $\DD$-\textit{fibration cohomologique complète,} je dis bien \underline{complète}, si

$a$) Pour tout $y \in \Ob\, Y$, $X_y \to X/y$ est une $\DD$-équivalence cohomologique (complète).

$b$) Pour tout $y' \to y$ dans $Y$, $X/y' \to X/y$ est une $\DD$-équivalence cohomologique (complète).

\printedpage{62}
% head: A. Grothendieck : Dérivateurs, Ch. I, Édition des Universités de Montpellier II et Paris VII
% foot: Version 8 juillet 2001

\msp{page 102}
$c$) Ces conditions restent satisfaites après tout changement de base $Y' \to Y$.

2°) Notion duale de fibration homologique complète.

\textsc{Exemple.} Si $f$ est \Hot-propre, les flèches des $a$) et $b$) sont des \Hot-équivalences, donc en vertu de \textbf{Der 3} $b'_1$) (qu’il faut admettre), ce sont des $\DD$-équivalences cohomologiques. Cela signifie que $a$) et $b$) sont valables, et $c$) aussi, car la notion de flèche \Hot-propre est stable par changement de base.

\textbf{Proposition 2.} \textit{Soient $f : X \to Y$ $\DD$-fibration cohomologique complète, $F \in \Ob\,\DD(X)$ localement constant sur $X$. Alors $f_*(F)$ est localement constant, sa formation commute à tout changement de base, et en particulier se calcule fibre par fibre.}\footnote[54]{Il y a une réciproque \textit{cf.} prop. 3$'$.}

\textbf{Corollaire 1.} \textit{Soit un diagramme cartésien}
\[
\begin{tikzcd}
X & X' \arrow[l, "q"'] \arrow[d, "f'"] \\
Y & Y' \arrow[l, "p"']
\arrow[from=1-1, to=2-1, "f"']
\end{tikzcd}
\quad .
\]
\textit{On suppose $f$ $\DD$-fibration cohomologique et $p$ une $\DD$-équivalence cohomologique. Alors $q$ est une $\DD$-équivalence cohomologique.}

\msp{page 103}
En effet, on a pour $F$ localement constant sur $X$
\[
H^*_{\DD}(X,F) \simeq H^*_{\DD}(Y, f_*F) \overset{\sim}{\longrightarrow} H^*_{\DD}(Y', p^*f_*F) \simeq H^*_{\DD}(Y', f'_*q^*F) \simeq H^*_{\DD}(X', q^*F)
\]
qed.

\textbf{Proposition 3.} \textit{La notion de $\DD$-fibration cohomologique complète est stable par changement de base} (trivial) \textit{et par composition.}

Pour la stabilité par composition, il suffit de noter le corollaire suivant de la proposition 2.

\textbf{Corollaire 2} (de prop. 2). \textit{Soit $f : X \to Y$. Pour que $f$ soit une $\DD$-fibration cohomologique complète, il faut et il suffit que pour tout $Y'' \overset{p}{\to} Y' \to Y$ sur $Y$, avec $p$ une $\DD$-équivalence cohomologique, la flèche correspondante $X'' = X \times_Y Y'' \longrightarrow X' = X \times_Y Y'$ soit aussi une $\DD$-équivalence cohomologique.}

C’est nécessaire par le corollaire 1, compte tenu que $f' : X' = X \times_Y Y' \to Y'$ est aussi une $\DD$-fibration cohomologique. C’est suffisant, car il suffit de l’appliquer aux cas $\{y\} \to Y/\{y\}$, $Y/y \to Y/y'$, qui sont des $\DD$-équivalences cohomologiques en
\msp{page 104}
vertu de \textbf{Der} 3 $b'_1$) (ou seulement \textbf{Der} 3 $a'$)).

\printedpage{63}
% head: A. Grothendieck : Dérivateurs, Ch. I, Édition des Universités de Montpellier II et Paris VII
% foot: Version 8 juillet 2001

\textbf{Proposition 4.} \textit{On suppose que} \textbf{Der} 3 a$'$) \textit{est satisfait} (\textit{i.e.} tout $Z$ dans \Dia{} qui est \Hot-asphérique, est complètement $\DD$-asphérique). \textit{Soit $f : X \to Y$ dans} \Dia. \textit{Les conditions suivantes sont équivalentes}

1°) \textit{$f$ est une $\DD$-fibration cohomologique \underline{complète}} (conditions $a$), $b$), $c$) de \texttt{[la]} définition 2, page 101).

2°) \textit{Pour tout $F \in$} \Ob\,$\DD(X)$\textit{, localement constant, $f_*(F)$ est localement constant, se calcule fibre par fibre; et ça reste vrai après tout changement de base $Y' \to Y$ (avec $F'$ localement constant sur $X' = X \times_Y Y'$). (Alors la formation de $f_*(F)$ commute à \underline{tout} changement de base).}

3°) \textit{Pour tout $Y'' \overset{p}{\to} Y'$ sur $Y$, si $p$ est une $\DD$-équivalence cohomologique, alors $X'' = X \times_Y Y'' \longrightarrow X' = X \times_Y Y'$ aussi.}

3$'$) \textit{Comme 3, en supposant que $Y''$ est la catégorie ponctuelle, et que $Y'$ a un élément final : il faut que $X'' \to X'$ soit une $\DD$-équivalence cohomologique.}

\msp{page 105}
\textbf{Définition 3.} On dit que $f : X \to Y$ dans \Dia{} est une $\DD$-\textit{fibration complète} si c’est une $\DD$-fibration cohomologique complète et une $\DD$-fibration homologique complète.

C’est stable par changement de base, par composition.

La définition 3 équivaut à ceci : pour tout $p : Y'' \to Y'$ sur $Y$, avec $Y''$ et $Y'$ \Hot-asphériques, le morphisme correspondant $X'' \to X'$ est \underline{à} \underline{la} \underline{fois} une $\DD$-équivalence cohomologique, et une $\DD$-équivalence homologique. (Et il suffit même de le supposer pour $Y'' = e$ et pour $Y'$ avec objet final ou objet initial.) On peut songer à renforcer encore cette condition en exigeant (sans regarder à gauche ni à droite) que $X'' \to X'$ soit une $\DD$-équivalence 1-complète. Mais il n’est pas clair que cette notion renforcée soit stable par composition – à moins de la renforcer encore, en exigeant que si $Y'' \to Y'$ est une $\DD$-équivalence 1-complète, alors $X'' \to X'$ l’est aussi. L’ennui avec cette définition renforcée, c’est que faute d’avoir une description cohomologique ou homologique de la notion de
\msp{page 106}
$\DD$-équivalence 1-complète, je ne me sens pas capable de donner des exemples intéressants (non totalement triviaux) de $\DD$-fibrations en ce sens très fort. Sauf les \Hot-fibrations triviales (\textit{i.e.} à fibres \Hot-asphériques) – mais ce n’est pas drôle!

Mais quand on a des produits $*$ associatifs et unitaires et des \underline{\Hom} internes correspondants dans les $\DD(X)_{\mathrm{loc.\ const.}}$, compatibles avec \texttt{[les]} images inverses, donc on aura sur $X$ si $F$ et $G$ sont localement constants dans $\DD(X)$, et si $\varepsilon_X$ désigne l’unité pour $*$ dans $\DD(X)$,
\[
\Hom(F,G) \simeq \Hom(\varepsilon_X * F, G) \simeq \Hom(\varepsilon_X, \underline{\Hom}(F,G)) \simeq \Hom_{\mathcal{A}}(\varepsilon_e, p_*\, \underline{\Hom}(F,G)) \ ,
\]
car $\varepsilon_X = p^*(\varepsilon_e)$, où $p : X \to e$, donc dans ce cas une flèche $f : X \to Y$ dans \Dia{} est une $\DD$-équivalence 1-complète si (et seulement si) c’est une $\DD$-équivalence cohomologique. Donc dans ce cas, si $f$ est une $\DD$-fibration cohomologique complète, c’est aussi une $\DD$-fibration homologique complète, donc aussi une $\DD$-fibration complète (définition 3). (Et même au sens soit-disant renforcé de tantôt : si $Y'' \to Y'$ sur $Y$ est une $\DD$-équivalence 1-complète ($\Leftrightarrow$ une $\DD$-équivalence cohomologique!) alors $X'' \to X'$ aussi. % [sic] soit-disant -> soi-disant, spelling; % [sic] "(Et même ... aussi." -> closing parenthesis missing, unbalanced parenthesis

% ---- from Der-ChI-234_p066-070.tex
\printedpage{64}
% head: A. Grothendieck : Dérivateurs, Ch. I, Édition des Universités de Montpellier II et Paris VII
% foot: Version 8 juillet 2001

\msp{page 107}

Il se pose la question de trouver des $f$ qui soient des $\DD$-fibrations cohomologiques, ou homologiques. Je ne sais si je serai capable d'en trouver, qui ne soient déjà des \HOT-fibrations cohomologiques, ou homologiques (cela revient maintenant au même\footnote[55]{Les \HOT-équivalences cohomologiques $=$ \HOT-équivalences homologiques $=$ \HOT-équivalences sans plus.} ), car pratiquement je ne sais guère dire qu'une flèche $Z' \to Z$ dans $\Cat$ est une $\DD$-équivalence cohomologique complète, que si cette flèche est une \HOT-équivalence (et en admettant l'axiome \textbf{Der} 3 $b'_1$), ou sa variante affaiblie cohomologique).

Je vais donc chercher les \HOT-fibrations, \textit{i.e.} les $f : X \to Y$ tels que pour tout $Y'' \to Y' \to Y$, avec $Y''$ et $Y'$ asphériques (sous-entendu : pour \HOT), $X'' \to X'$ est une \Hot-équivalence. (Il suffit même de supposer $Y''$ réduit à un point, $Y'$ avec objet final $\ldots$ )

Supposons par exemple $f$ propre (\textit{i.e.} \HOT-propre), donc la condition $a$) de la définition 2, p. 101 satisfaite et de même après tout changement de base. Suffit-il de vérifier de plus la condition $b$), \textit{i.e.} que
\[ X/y \longrightarrow X/y' \quad , \qquad \text{pour } y \to y' \text{ dans } Y \ , \]
sont des \Hot-équivalences? (J'entends en restant sur $Y$, sans supposer aussi d'avoir la même chose après changement de base).
\msp{page 108}
Mais cela reste-il vrai après tout changement de base? % [sic] reste-il -> reste-t-il, forme interrogative

La situation semble ici meilleur pour $f$ lisse (au lieu de propre), en interprétant les conditions $a$) et $b$) sous forme homologique, comme signifiant que $f_!(e_X)$ est un objet \underline{localement} \underline{constant} de $\mathsf{HOT}(Y)$, qui se calcule fibre par fibre. Comme $f$ lisse implique que $f_!$ commute à tout changement de base, cette hypothèse sur $f$ (que $f$ soit une \HOT-fibration cohomologique, peut-être pas complète) est stable par changement de base, donc implique que c'est une \HOT-fibration cohomologique complète. Mais cet argument utilise les propriétés générales des dérivateurs pour \HOT, que je n'ai pas encore établies - surtout pas dans le domaine $\Cat$ tout entier! Donc il est sain de me méfier. Je préfère tâter la difficulté en restant d'abord dans le cadre du formalisme \underline{cohomologique} (non commutatif), dont j'ai une certaine maîtrise. Je suppose donc $f$ propre à nouveau, et que les % [sic] meilleur -> meilleure, accord avec « situation »
\[ X_y \longrightarrow X/y \quad , \qquad X/y \longrightarrow X/y' \]
\msp{page 109}
soient des \HOT-équivalences, je veux prouver la même chose après tout changement de base, au moins que dans la situation
\[
\begin{tikzcd}
X \arrow[r, leftarrow, "q"] \arrow[d, "f"'] & X' \arrow[r, leftarrow, "q'"] \arrow[d, "f'"'] & X'' \arrow[d, "f''"] \\
Y \arrow[r, leftarrow, "p"'] & Y' \arrow[r, leftarrow, "p'"'] & Y''
\end{tikzcd}
\]
\printedpage{65}
% head: A. Grothendieck : Dérivateurs, Ch. I, Édition des Universités de Montpellier II et Paris VII
% foot: Version 8 juillet 2001
si $Y'$, $Y''$ sont asphériques, alors $q'$ est une \Hot-équivalence. Si je veux l'établir par critère cohomologique, je dois commencer par prendre sur $X'$ un faisceau d'\textit{ensembles} localement constant $F$, et prouver que $\Gamma(X',F) \overset{\sim}{\to} \Gamma(X'',{q'}^{*}F)$. Comme ce faisceau ne provient pas de $X$, je suis ennuyé pour appliquer l'hypothèse sur $f$, qui concerne à priori des faisceaux d'ensembles constants sur $X$, ou sur les $X/y$. Mais on devrait avoir ceci :

\textbf{Lemme.} \textit{Soit $f : X'' \to X'$. Alors $f$ est une} \HOT-\textit{équivalence} i.e.
\[ H^{*}_{\mathsf{HOT}}(X', a_{X'}) \xrightarrow{\ \sim\ } H^{*}_{\mathsf{HOT}}(X'', a_{X''}) \]
\textit{est} \texttt{[un]} \textit{isomorphisme pour tout $a \in \mathsf{Hot} = \mathsf{HOT}(e)$, si et seulement si $f$ est une} \Hot-\textit{équivalence,}\footnote[56]{Attention ici \texttt{[mot illisible]} la distinction entre \HOT-équivalence et \Hot-équivalence.} \textit{et alors $f$ est même une} \HOT-\textit{équivalence} $2$-\textit{complète} (j'ai besoin seulement de : \HOT-équivalence cohomologique et homologique complète) i.e. \textit{pour tout $\xi \in \DD(X')$ \underline{localement} \underline{constant} sur $X'$, on a}\footnote[57]{Est-ce vrai aussi pour les dérivateurs $\DD_i$?}
\[ H^{*}_{\mathsf{HOT}}(X',\xi) \xrightarrow{\ \sim\ } H^{*}_{\mathsf{HOT}}(X'',{q'}^{*}(\xi)) \quad . \]
\msp{page 110}

La première assertion doit résulter des formules
\[ H^{*}_{\mathsf{HOT}}(Z,a_Z) \simeq \underline{\mathsf{Hom}}_{\mathsf{Hot}}(\mathsf{hot}(Z),a) \tag*{$(*)$} \]
\[ \text{(et aussi} \qquad\qquad H_{*}^{\mathsf{HOT}}(Z,a_Z) \simeq \mathsf{hot}(Z) \times a \qquad\qquad ) \]
donc si
\[ \mathsf{hot}(X'') \xrightarrow{\ \sim\ } \mathsf{hot}(X') \tag*{$(**)$} \]
est \texttt{[un]} isomorphisme, alors
\[ H^{*}_{\mathsf{HOT}}(X', a_{X'}) \longrightarrow H^{*}_{\mathsf{HOT}}(X'', a_{X''}) \]
est \texttt{[un]} isomorphisme, puisque c'est déduit de $(*)$ en prenant $\underline{\mathsf{Hom}}_{\mathsf{Hot}}(\,.\,,a)$. Inversement, si $(*)$ est toujours \texttt{[un]} isomorphisme, exprimant les deux membres \texttt{[mot illisible]} de $\underline{\mathsf{Hom}}_{\mathsf{Hot}}(\,.\,,a)$ et prenant les $\pi_0$, on trouve que
\[ \mathsf{Hom}_{\mathsf{Hot}}(\mathsf{hot}(X'),a) \xrightarrow{\ \sim\ } \mathsf{Hom}_{\mathsf{Hot}}(\mathsf{hot}(X''),a) \]
est un isomorphisme pour tout $a$, donc $(**)$ est un isomorphisme.

Mais le fait que ça implique que $f''$ soit une \HOT-équivalence cohomologique complète, \textit{i.e.} que la relation $(*)$ reste vraie pour un $\xi \in \DD(X')$ \underline{localement} \underline{constant} sur $X'$, me paraît une propriété délicate du dérivateur \HOT, sur laquelle il me faudra revenir.

Donc le lemme est loin d'être établi, faute d'avoir seulement construit le dérivateur \HOT, sans même parler d'en établir les propriétés spéciales délicates! Mais je n'ai aucune raison de douter du lemme, il fait partie de la kyrielle des intuitions de base que je voudrais fonder solidement.

\printedpage{66}
% head: A. Grothendieck : Dérivateurs, Ch. I, Édition des Universités de Montpellier II et Paris VII
% foot: Version 8 juillet 2001

\msp{page 111}

Le lemme implique qu'une $\DD = \mathsf{HOT}$-fibration (au sens de la définition 2 page 81, revue et corrigée à la page 96-97) est en fait déjà une \HOT-fibration cohomologique et homologique complète. Or on a vu que si $f$ est \underline{propre}, $f$ est une $\DD$-fibration si et seulement si les conditions précédentes \underline{sur $Y$} (sans changement de base) sont satisfaites. Et dualement pour $f$ lisse.

\begin{flushright}
11.11.90
\end{flushright}

Je sens le besoin de revoir les notions de $\DD$-fibrations en tous genres (p. 76 ff). D'abord les diverses variantes utiles de l'axiome \textbf{Der} 3 sur les dérivateurs :

\textbf{Der} 3 $a$)\footnote[58]{\textit{cf.} aussi la série \textbf{Der 3} $c \ldots$ , page 124 ff.} Si $Z \in \mathsf{Dia}$ est \Hot-asphérique, il est $\DD$-asphérique \textit{i.e.} $\mathcal{A} = \DD(e) \to \DD(Z)$ est pleinement fidèle. % Doubt resolved: attached PDF leaf 30 (printed p.66), footnote 58 has bold Der 3.

$a'$) Comme $a$), avec en plus la caractérisation de l'image essentielle de $\mathcal{A}$ dans $\DD(Z)$, comme formée des $\DD$-coefficients localements constants. % [sic] localements -> localement, adverbe invariable

\textbf{Der} 3 $b$) Si $f : Z' \to Z$ dans \Dia{} est une \Hot-équivalence, alors c'est une $\DD$-équivalence, \textit{i.e.} le foncteur $f^{*} : \DD(Z) \to \DD(Z')$ induit une équivalence entre les sous-catégories formées des coefficients constants (ou encore un foncteur pleinement fidèle $\DD(Z)_{\text{loc. const.}} \longrightarrow \DD(Z')$).

$b'_1$) Si $f$ comme ci-dessus, alors $f$ est une $\DD$-équivalence 1-complète, \textit{i.e.} le foncteur $f^{*} : \DD(Z) \to \DD(Z')$ induit un foncteur pleinement fidèle
\msp{page 112}
entre les sous-catégories pleines formées des objets localement constants.

$b'_2$) Comme $b'_1$), mais avec $\DD$-équivalence 2-\underline{complète} au lieu de 1-complète, \textit{i.e.} $f^{*}$ induit une équivalence entre les sous-catégories formées des coefficients localement constants.

$b''_{\mathrm{coh}}$) Si $f : Z' \to Z$ dans \Dia{} est une \Hot-équivalence, alors c'est une $\DD$-équivalence cohomologique complète, \textit{i.e.} pour tout $F \in \mathsf{Ob}\ \DD(Z)$ qui est \underline{localement} \underline{constant}, on a
\[ H^{*}_{\DD}(Z,F) \xrightarrow{\ \sim\ } H^{*}_{\DD}(Z',f^{*}F) \quad . \]

$b''_{\mathrm{hom}}$) Si $f : Z' \to Z$ dans \Dia{} est une \Hot-équivalence, c'est une $\DD$-équivalence homologique complète, \textit{i.e.} pour tout $F \in \mathsf{Ob}\ \DD(Z)$ qui est \underline{localement} \underline{constant}, on a\footnote[59]{\textit{cf.} variante $\pm$ spéciale pour \HOT{} pages 124-125.}
\[ H_{*}^{\DD}(Z',f^{*}F) \xrightarrow{\ \sim\ } H_{*}^{\DD}(Z,F) \quad . \]

\printedpage{67}
% head: A. Grothendieck : Dérivateurs, Ch. I, Édition des Universités de Montpellier II et Paris VII
% foot: Version 8 juillet 2001

On a le diagramme d'implications tautologiques pour les variantes de \textbf{Der} 3 :
\[
\begin{tikzcd}
 & b'_2 \arrow[d, Rightarrow] \arrow[drr, Rightarrow] & & & \\
 & b'_1 \arrow[dl, Rightarrow] \arrow[dr, Rightarrow] \arrow[dd, Rightarrow] & & a' \arrow[ddd, Rightarrow] & \\
b''_{\mathrm{coh}} \arrow[dr, Rightarrow] & & b''_{\mathrm{hom}} \arrow[dl, Rightarrow] & & \\
 & b \arrow[drr, Rightarrow] & & & \\
 & & & a & .
\end{tikzcd}
\]
Dans tous les cas à ma connaissance, les dérivateurs satisfont la condition la plus forte \textbf{Der} 3 $b'_2$.
\msp{page 113}

Je rappelle aussi la :

\textbf{Définition 1 :} Soit $f : X \to Y$ dans \Dia. On dit que $f$ est une $\DD$-\textit{équivalence cohomologique} (sous-entendu : “\textit{complète}”) si pour tout $\DD$-coefficient $F$ \underline{localement} \underline{constant} sur $Y$,
\[ H^{*}_{\DD}(Y,F) \xrightarrow{\ \sim\ } H^{*}_{\DD}(X,f^{*}F) \quad ; \]
on dit que $f$ est une $\DD$-\textit{équivalence homologique} (“\textit{complète}”) si pour tout $F$ comme \texttt{[ci-]}dessus,
\[ H_{*}^{\DD}(X,f^{*}F) \xrightarrow{\ \sim\ } H_{*}^{\DD}(Y,F) \quad . \]
On dit que $f$ est une $\DD$-\textit{équivalence} 1-\textit{complète} (resp. 2-\textit{complète}) si le foncteur
\[ f^{*} : \DD(Y)_{\text{loc. const.}} \longrightarrow \DD(X)_{\text{loc. const.}} \]
est pleinement fidèle (resp. une équivalence).
\begingroup\small
\[
\begin{tikzcd}[column sep=small]
 & \text{$\DD$-équivalence 2-complète} \arrow[d, Rightarrow] & & \\
 & \text{$\DD$-équivalence 1-complète} \arrow[dl, Rightarrow] \arrow[dr, Rightarrow] & & \\
\text{$\DD$-équivalence cohomologique} \arrow[dr, Rightarrow] & & \text{$\DD$-équivalence homologique} \arrow[dl, Rightarrow] & \\
 & \text{$\DD$-équivalence} & & .
\end{tikzcd}
\]
\endgroup
Si dans les sous-catégories $\DD(X)_{\text{loc. const.}}$ il y a une loi $*$ (associative, unitaire, commutative) et des \underline{\textsf{Hom}} internes associés, préservés par les images inverses, alors les $\DD$-équivalences cohomologiques sont déjà 1-complètes et sont donc des $\DD$-équivalences homologiques, donc dans ce cas particulier
\printedpage{68}
% head: A. Grothendieck : Dérivateurs, Ch. I, Édition des Universités de Montpellier II et Paris VII
% foot: Version 8 juillet 2001
\msp{page 114 vide}
\msp{page 115}
\[
\begin{tikzcd}
\text{$\DD$-équivalence 2-complète} \arrow[d, Rightarrow] & & \\
\text{$\DD$-équivalence 1-complète} \arrow[r, Leftrightarrow] & \text{$\DD$-équivalence cohomologique} \arrow[d, Rightarrow] & \\
 & \text{$\DD$-équivalence homologique} \arrow[d, Rightarrow] & \\
 & \text{$\DD$-équivalence} & .
\end{tikzcd}
\]
Dans $\DD = \mathsf{HOT}$, toutes ces notions devraient être équivalentes, et équivalentes à la condition encore plus faible en apparence, que $f : X \to Y$ soit une \Hot-équivalence. Mais je ne connais que \HOT{} comme dérivateur où la $\DD$-équivalence d'une flèche $f : X \to Y$ dans \Dia{} implique déjà la $\DD$-équivalence cohomologique (complète), ou la $\DD$-équivalence homologique (complète), voire même la $\DD$-équivalence 1-complète. (Je reviendrai là-dessus tantôt.)

\textbf{Proposition 1.} \textit{Soit $f : X \to Y$ dans} \Dia. \textit{Conditions équivalentes :}

a) \textit{Pour tout $y \in \mathsf{Ob}\, Y$, l'inclusion $X_y \to X/y$ est une $\DD$-équivalence, et pour toute flèche $y' \to y$ dans $Y$, $X/y' \to X/y$ est une $\DD$-équivalence.}

b) \textit{Pour tout coefficient \underline{constant} $F$ sur $Y$, $f_{*}f^{*}F$ est localement constant, et se calcule “fibre par fibre”.}

(C'est la prop. 1, p. 76).

\textbf{Proposition 1$'$.} Dual.

\msp{page 116}

\textbf{Définition 2.} Sous les conditions équivalentes de la proposition 1, on dit que $f$ est une $\DD$-\textit{fibration cohomologique faible.} Sous les conditions duales, que $f$ est une $\DD$-\textit{fibration homologique faible.}

C'est ce que dans la définition 1, p. 77 j'appelais “$\DD$-fibration cohomologique” ou “$\DD$-fibration à gauche”, et dualement. Maintenant la notion me paraît trop faible pour mériter ce nom, que je réserve pour le cas “complet”.

\textbf{Proposition 2.} \textit{Les $\DD$-fibrations cohomologiques faibles sont stables par changements de base lisses, les $\DD$-fibrations homologiques faibles stables par changements de base propres.}

Je devrais construire un exemple où $f$ est une $\DD$-fibration cohomologique faible, et ne l'est plus après changement de base.

% ---- from Der-ChI-234_p071-075.tex
\printedpage{69}
% head: A. Grothendieck : Dérivateurs, Ch. I, Édition des Universités de Montpellier II et Paris VII

\textbf{Proposition 3.} \textit{Soit}
\[
\begin{tikzcd}
X \arrow[rr, "f"] \arrow[dr, "p"'] & & X' \arrow[dl, "p'"] \\
& S &
\end{tikzcd}
\]
\textit{dans \Dia, avec $p$ et $p'$ des $\D$-fibrations cohomologiques faibles} (resp. \textit{des $\D$-fibrations homologiques faibles}), \textit{soit $S_0 \subset S$ tel que $\pi_0(S_0) \to \pi_0(S)$ soit surjective. Pour que $f$ soit une $\D$-équivalence, il suffit que les $f_s : X_s \to X'_s$, pour $s \in S_0$, le soient.}

\textbf{Proposition 4.} \textit{Si $f : X \to Y$ dans \Dia{} est $\D$-propre} (resp. $\D$-\textit{lisse}), \textit{alors $f$ est une $\D$-fibration cohomologique} (resp. \textit{homologique}) \textit{faible si et seulement si pour tout $F$ constant sur $Y$, $f_*f^*F$} (resp. $f_!f^*F$) \textit{est localement constant sur $Y$. Et s'il en est ainsi, pour tout $Y' \to Y$ dans \Dia, $f' : X' = X \times_Y Y' \to Y'$ est encore une $\D$-fibration cohomologique} (resp. \textit{homologique}) \textit{faible.}

Immédiat sur le critère \textit{b}) de la proposition 1.

\msp{page 117}\textbf{Proposition 5.} \textit{Soit $f : X \to Y$ dans \Dia. Considérons les conditions suivantes :}

a) \textit{Pour tout $Y'' \to Y'$ sur $Y$ dans \Dia, si $Y''$ et $Y'$ sont \Hot-asphériques, alors}
\[
X'' = X \times_Y Y'' \longrightarrow X' = X \times_Y Y'
\]
\textit{est une $\D$-équivalence.}
\[
\begin{tikzcd}
X & X' \arrow[l] & X'' \arrow[l] \\
Y \arrow[from=u] & Y' \arrow[l] \arrow[from=u] & Y'' \arrow[l, "p"] \arrow[from=u]
\end{tikzcd}
\]

a$'$) \textit{Comme} a), \textit{mais en supposant $p : Y'' \to Y'$ une \Hot-équivalence (au lieu de supposer $Y'$, $Y''$ \Hot-asphériques).}

a$_1$) \textit{Comme} a), \textit{mais en supposant (au lieu de l'asphéricité) que $Y'$ a un objet final, et $Y''$ est la catégorie unité.}

a$_2$) \textit{Comme} a), \textit{mais en supposant que $Y'$ a un objet cofinal, et $Y''$ est la catégorie unité.}

b$_1$) \textit{$f$ est une $\D$-fibration cohomologique faible, et le reste après tout changement de base.}

b$_2$) \textit{$f$ est une $\D$-fibration homologique faible, et le reste après tout changement de base.}

\printedpage{70}
% head: A. Grothendieck : Dérivateurs, Ch. I, Édition des Universités de Montpellier II et Paris VII
\textit{On a alors le diagramme d'implications tautologiques}
\[
\begin{tikzcd}[column sep=small]
& & & a' \arrow[d, Rightarrow] & & & & \\
& & & a \arrow[dlll, Rightarrow] \arrow[drrr, Rightarrow] & & & & \\
a_1 & \Longleftrightarrow & b_1 & & b_2 & \Longleftrightarrow & a_2 & .
\end{tikzcd}
\]
\textit{Lorsque $\D$ satisfait à l'axiome} \textbf{Der} 3 a$'$ (p. 111), \textit{alors toutes ces conditions, sauf} a$'$), \textit{sont équivalentes. Lorsque $\D$ satisfait à l'axiome} \textbf{Der} 3 b$'_1$ \textit{(ou seulement} \textbf{Der} 3 b$''_{\mathrm{coh}}$ ou b$''_{\mathrm{hom}}$\textit{),} \textit{alors toutes les conditions (y inclus} a$'$\textit{) sont équivalentes.}

Pour simplifier, on va par la suite supposer que l'axiome \textbf{Der} 3 $a'$) est satisfait.

\textbf{Définition 3.} Sous les conditions équivalentes $a$, $a_1$, $b_1$, $a_2$, $b_2$ de la proposition 4, on dit que $f$ est une $\D$-\textit{fibration}. % [sic] proposition 4 -> proposition 5, the conditions a, a_1, b_1, a_2, b_2 are those of proposition 5

\msp{page 118}\textbf{Proposition 6.} \textit{Supposons que $f : X \to Y$ dans \Dia{} soit $\D$-propre} (resp. $\D$-\textit{lisse}). \textit{Conditions équivalentes :}

a) \textit{$f$ est une $\D$-fibration.}

b) \textit{$f$ est une $\D$-fibration cohomologique} (resp. \textit{homologique}) \textit{faible.}

c) \textit{Pour tout $y' \to y$ dans $Y$, $X/y' \to X/y$} (resp. $y\backslash X \to y'\backslash X$) \textit{est une $\D$-équivalence.}

d) \textit{Pour tout $F$ constant sur $Y$, $f_*f^*F$} (resp. $f_!f^*F$) \textit{est localement constant sur $Y$.}

d$'$) \textit{Pour tout $F$ localement constant sur $Y$, $f_*f^*F$} (resp. $f_!f^*F$) \textit{est localement constant sur $Y$.}

Cela résulte formellement de Prop. 4 et Prop. 5.

\textbf{Corollaire.} \textit{Supposons $f : X \to Y$} \Cat\textit{-fibrant} (resp. \Cat\textit{-cofibrant}). \textit{Pour que $f$ soit une $\D$-fibration, il faut et il suffit que pour tout $u : y \to y'$ dans $Y$, $u^* : X_{y'} \to X_y$} (resp. $u_* : X_y \to X_{y'}$) \textit{soit une $\D$-équivalence.}

Il est immédiat que la notion de $\D$-fibration est stable par changement de base, mais il est douteux qu'elle le soit par composition.\footnote[60]{N.B. Si les $\D$-équivalences sont déjà des \Hot-équivalences, alors la notion coïncide avec celle de \Hot-fibration, et elle est stable par composition. Également, si les $\D$-équivalences sont au sens des $\D$-équivalences cohomologiques, ou homologiques, alors la notion coïncide avec celle des $\D$-fibrations cohomologiques (resp. homologiques) ci-dessous, et est également stable par composition.} (Ici encore, j'ai omis de faire un contre-exemple, avec des dérivateurs $\D = \sfup{Der}(k\text{-}\sfup{Mod})$ disons.) Autre lacune de la notion : il n'est pas clair, si $f$ est une $\D$-fibration et $F$ un $\D$-coefficient sur $Y$, si
\[
f_*f^*F \qquad \text{et} \qquad f_!f^*F
\]
\printedpage{71}
% head: A. Grothendieck : Dérivateurs, Ch. I, Édition des Universités de Montpellier II et Paris VII
se calculent “fibre par fibre”, comme on s'attendrait dans le cas d'un “fibré”.

Par contre, voici un résultat sympa :

\msp{page 119}\textbf{Proposition 7.} \textit{Soit $f : X \to Y$ dans \Dia. Conditions équivalentes :}

a) \textit{$f$ est une $\D$-équivalence “universelle”,} i.e. \textit{pour tout $Y' \to Y$ dans \Dia, $f' : X' = X \times_Y Y' \to Y'$ est une $\D$-équivalence.}

b) \textit{$f$ est une $\D$-fibration à fibres $\D$-asphériques.}

\textbf{Définition 4.} Sous les conditions équivalentes de la proposition 7, on dit que $f$ est une $\D$-fibration \textit{triviale}.

Mais attention, $f$ peut-être une $\D$-fibration et une $\D$-équivalence, sans que la $\D$-fibration soit triviale \textit{i.e.} sans que les fibres soient $\D$-asphériques, cf. p. 87 et suivantes. % [sic] peut-être -> peut être, verb "pouvoir" followed by infinitive
Et également (devrait remonter) même si $\D$ satisfait \textbf{Der} 3 $b_2$ (la plus forte des variantes de \textbf{Der} 3), donc quand \underline{toutes} les conditions de la proposition 5, y inclus $a'$), sont équivalentes, celles-ci (\textit{i.e.} que $f$ soit une $\D$-fibration) n'impliquent pourtant pas, en général, que dans la condition $a'$) on puisse se borner à $p : Y'' \to Y'$ une $\D$-équivalence (au lieu d'une \Hot-équivalence). En d'autres termes, si on a un diagramme cartésien dans \Dia
\[
\begin{tikzcd}
X \arrow[d, "f"'] & X' \arrow[l, "q"'] \arrow[d, "f'"] \\
Y & Y' \arrow[l, "p"] & ,
\end{tikzcd}
\]
avec $f$ une $\D$-fibration, et $p$ une $\D$-équivalence, il ne s'ensuit pas forcément que $q$ soit une $\D$-équivalence. (Mais ce sera pourtant le cas, “bien sûr”, si $\D = \HOT$.)

\msp{page 120}Par contre, on a :

\textbf{Proposition 8.} \textit{Soit $f : X \to Y$ dans \Dia. Conditions équivalentes :}

a) \textit{$f$ est une $\D$-fibration.}

b) \textit{Pour tout $p : Y'' \to Y'$ sur $Y$, si $p$ est une $\D$-équivalence} 1-\textit{complète} (déf. 1, p. 113), \textit{alors $q : X'' \to X'$ est une $\D$-équivalence.}
\[
\begin{tikzcd}
X \arrow[d, "f"'] & X' \arrow[l] \arrow[d, "f'"'] & X'' \arrow[l, "q"'] \arrow[d, "f''"] \\
Y & Y' \arrow[l] & Y'' \arrow[l, "p"]
\end{tikzcd}
\]
\textbf{Corollaire.} \textit{Supposons que dans \Dia, toute $\D$-équivalence soit} 1-\textit{complète.} (N.B. Je ne connais que le dérivateur \HOT{} pour satisfaire à cette condition, qui paraît quasiment caractéristique de \HOT.) \textit{Alors $f : X \to Y$ est une $\D$-fibration si et seulement si pour tout $p : Y'' \to Y'$ sur $Y$, si $p$ est une $\D$-équivalence, alors $q : X'' \to X'$ aussi. Dans ce cas exceptionnel parmi les dérivateurs, la notion de $\D$-fibration est stable par composition} (ce qui est tautologique sur le critère que je viens de donner.)

\printedpage{72}
% head: A. Grothendieck : Dérivateurs, Ch. I, Édition des Universités de Montpellier II et Paris VII

On va renforcer la notion de $\D$-fibration, pour la rendre plus souple, et notamment stable par composition.

\textbf{Proposition 9.} \textit{Soit $f : X \to Y$ dans \Dia. Conditions équivalentes :}

a) \textit{Pour tout $y \in Y$, $X_y \to X/y$ est une $\D$-\underline{équivalence} \underline{cohomologique}} (déf. 1, p. 113), \textit{et de même pour $X/y \to X/y'$ pour toute flèche $y \to y'$ dans $Y$. Et ces conditions sur $f$ restent valables après tout changement de base} (N.B. il suffit que la première de ces conditions le reste).

\msp{page 121}b) \textit{pour tout $F \in \Ob\ \D(X)$ \underline{localement} \underline{constant}, $f_*F$ est localement constant sur $Y$ et se calcule “fibre par fibre” -- et itou après tout changement de base.} % Doubt resolved: attached PDF leaf 36 (printed p.72) has an en dash, encoded --.

c) \textit{Pour tout $p : Y'' \to Y'$ sur $Y$, si $p$ est une $\D$-équivalence cohomologique, alors il en est de même de $q : X'' \to X'$.}

c$'$) \textit{Comme} c, \textit{mais en supposant comme hypothèse sur $p$ que $Y''$ soit la catégorie finale, et $Y'$ ait un objet final.}

\textsc{Démonstration}. Dans l'ordre ($\pm$ tautologique)
\[
\mathrm{a} \Longrightarrow \mathrm{b} \Longrightarrow \mathrm{c} \Longrightarrow \mathrm{c}' \Longrightarrow \mathrm{a} \quad .
\]

\textbf{Définition 5.} Sous les conditions équivalentes de la proposition 9, on dit que $f$ est une $\D$-\textit{fibration cohomologique}. Sous les conditions duales, on dit que $f$ est une $\D$-\textit{fibration homologique}. On dit que $f$ est une $\D$-\textit{fibration forte}, si c'est à la fois une $\D$-fibration cohomologique et une $\D$-fibration homologique.

Ces notions renforcent la notion de $\D$-fibration. Elles ont l'avantage de plus d'être \underline{stables} \underline{par} \underline{composition}, non seulement par changement de base.

\msp{page 122}La notion de $\D$-fibration forte a été choisie pour coiffer les deux notions duales de $\D$-fibration cohomologique et homologique (à défaut d'une meilleur notion) tout en restant stable par composition. On pourrait exiger que pour tout % [sic] meilleur -> meilleure, agreement with "notion"
\[
Y'' \longrightarrow Y'
\]
sur $Y$, avec $Y'$, $Y''$ \Hot-asphériques, $q : X'' \to X'$ soit une $\D$-équivalence 1-complète (voire même 2-complète), appeler ça les $\D$-\textit{fibrations} 1-\textit{complètes} (resp. 2-\textit{complètes}). Ça renforce de façon assez élégante la notion de $\D$-fibration forte, mais a le grand tort de ne pas être stable par composition (que je sache!), sauf dans le cas particulier où on a des \underline{\Hom} internes sympas (cf. bas de la page 113), auquel cas la notion de $\D$-équivalence implique déjà celle de $\D$-équivalence 1-complète,
\printedpage{73}
% head: A. Grothendieck : Dérivateurs, Ch. I, Édition des Universités de Montpellier II et Paris VII
\msp{page 123}donc aussi celle de $\D$-équivalence homologique. Donc dans ce cas
\[
\begin{tikzcd}
\text{$\D$-fibration cohomologique} \arrow[r, Leftrightarrow] \arrow[d, Rightarrow] & \text{$\D$-fibration 1-complète} \\
\text{$\D$-fibration homologique} \quad , &
\end{tikzcd}
\]
et il n'y a pas à chercher plus loin que les $\D$-fibrations cohomologiques.

On pourrait songer, il est vrai, à définir (dans le cas général) une notion plus forte, en exigeant que pour tout $p : Y'' \to Y'$ sur $Y$ qui soit une $\D$-équivalence 1-complète, $q : X'' \to X'$ soit aussi une $\D$-équivalence 1-complète. Cette fois-ci, la notion est stable par composition, mais elle me paraît néanmoins inutilisable (en dehors du cas qu'on vient de dire, où elle coïncide avec la notion de $\D$-fibration cohomologique, ou encore de $\D$-fibration forte), car n'étant pas en mesure de décrire les $\D$-équivalences 1-complètes, je suis aussi incapable de donner des exemples non triviaux de $\D$-fibrations “strictement 1-complètes”.

Pour vérifier qu'une $f : X \to Y$ dans \Dia{} est une $\D$-fibration cohomologique, il faut
\msp{page 124}savoir vérifier que des flèches $q : Z' \to Z$ dans \Dia{} sont des $\D$-équivalences cohomologiques, donc que $H^*_{\D}(Z,F) \overset{\sim}{\rightarrow} H^*_{\D}(Z',q^*(F))$ est \texttt{[un]} isomorphisme pour tout $F \in \Ob\ \D(Z)$ qui est localement constant sur $Z$. Or je ne sais vérifier un énoncé de cette nature (en dehors du cas où je sais d'avance que $Z$ est complètement $\D$-asphérique, \textit{i.e.} que les coefficients localement constants sur $Z$ sont en fait constant) \underline{que} si $q$ est déjà une \Hot-équivalence. % [sic] constant) -> constants), agreement with "coefficients"
C'est dire que je ne sais exhiber des $\D$-fibrations cohomologiques, que quand ce sont déjà des \HOT-fibrations cohomologiques. Je veux, dans la suite, donner des critères utilisables pour des \HOT-fibrations. Mais on va axiomatiser ces développements, en gardant un $\D$ général, sur lequel je vais faire une hypothèse si draconienne que pratiquement, je ne connais que le dérivateur \HOT{} qui y satisfait. Il s'agit de $c_1$ ci dessous.

\textbf{Der} 3 $c_1$) Toute $\D$-équivalence est une $\D$-équivalence cohomologique.

$c_2$) Toute $\D$-équivalence est une $\D$-équivalence homologique.

\msp{page 125}$c'$) Toute $\D$-équivalence est 1-complète.

$c''$) Toute $\D$-équivalence est 2-complète.

N.B. On a
\[
\begin{tikzcd}
& c'' \arrow[d, Rightarrow] & \\
& c' \arrow[dl, Rightarrow] \arrow[dr, Rightarrow] & \\
c_1 & & c_2
\end{tikzcd}
\]

% ---- from Der-ChI-234_p076-076.tex
\printedpage{74}
% head: A. Grothendieck : Dérivateurs, Ch. I, Édition des Universités de Montpellier II et Paris VII
% foot: Version 8 juillet 2001
tautologiquement. Je ne connais que \HOT{} pour satisfaire l'ensemble des deux plus faibles conditions $c_1$, $c_2$. Et en fait, sous réserve d'écrire la théorie du dérivateur \HOT, il doit être vrai qu'elle satisfait même la forme la plus forte $c''$. (Et les \HOT-équivalences doivent être simplement les \Hot-équivalences.) (Cf. le “lemme” p.~109.)

La proposition suivante est une tautologie.

\noindent\textbf{Proposition 10.} \textit{Supposons que $\D$ satisfasse la condition \textup{\textbf{Der} 3 $\mathrm{c}_1$ (p.~124),} par exemple $\D = \text{\HOT}$. Alors la notion de $\D$-fibration coïncide avec celle de $\D$-fibration cohomologique. Elle est donc stable par composition. Pour que $f : X \to Y$ soit une $\D$-fibration, il faut et il suffit que pour $p : Y'' \to Y'$ sur $Y$, si $p$ est une $\D$-équivalence, il en est de même de $q : X'' \to X'$.}

\noindent\textbf{Corollaire 1.} \textit{Si $f : X \to Y$ dans \textup{\Dia} est $\D$-propre, alors $f$ est une $\D$-fibration cohomologique si et seulement si pour toute flèche $y' \to y$ dans $Y$, $X/y' \to X/y$ est une $\D$-équivalence \textup{(i.e.} si et seulement si $f$ est une $\D$-fibration cohomologique faible\textup{).}}

Appliquer le critère $c$) de \texttt{[la]} prop.~6.

\msp{page 126}

\noindent\textbf{Corollaire 2.} \textit{Supposons que $\D$ satisfasse \textup{\textbf{Der} 3 $(\mathrm{c}_1, \mathrm{c}_2)$ (par exemple $\D = \text{\HOT}$).} Alors pour une flèche $f : X \to Y$ dans \textup{\Dia}, il revient au même que $f$ soit une $\D$-fibration, une $\D$-fibration cohomologique, ou une $\D$-fibration homologique. Si $f$ est propre \textup{(resp.} lisse\textup{),} il revient au même de dire que $f$ est une $\D$-fibration cohomologique \textup{(resp.} homologique\textup{),} faible.}

On voit donc que dans $\Cat$, on doit avoir une seule notion raisonnable de \HOT-fibration, qui coïncidera avec celle de \HOT-fibration cohomologique ou homologique. Mais pour que cette notion soit dès à présent disponible, il faudrait que le dérivateur \HOT{} soit déjà construit sur le domaine $\text{\Dia} = \Cat$ tout entier. Or ce n'est pas chose faite encore (et il n'a donc pas pu être question de prouver les propriétés de \HOT{} dont il vient d'être question, disons la plus forte \textbf{Der} 3 $c''$).

Néanmoins, je voudrais développer de façon directe (“élémentaire”) la notion de \Hot-fibration, sans me réferer à la théorie du dérivateur \HOT, et aux développements théoriques précédents. % [sic] réferer -> référer, missing accent on the second e
\end{document}
